\paragraph{Green 核的谱代数：Gram 分解、Chebyshev 识别与重整化常数}
本段补充说明锯齿核 \(K_k(i,j)=\min(i,j)\) 的若干\emph{整数/谱/行列式}刚性后果。
这些结论把先前出现的 \(4\sin^2(\cdot)\) 形式从“半角恒等式”提升为可审计的\emph{Gram 归一化 + Chebyshev 整系数}的闭合约束链。

\begin{lemma}[Gram 分解与行列式归一化]\label{lem:pom-Kk-gram-det}\leanverified{paper\_pom\_kk\_gram\_det}
令 \(K_k\in\ZZ^{k\times k}\) 由 \(K_k(i,j)=\min(i,j)\) 定义。
取单位下三角矩阵 \(B_k\in\ZZ^{k\times k}\) 满足
\[
(B_k)_{ij}:=\mathbf{1}_{\{j\le i\}}.
\]
则有严格 Gram 分解
\[
K_k=B_kB_k^{\mathsf T},
\qquad
\det(B_k)=1,
\qquad
\det(K_k)=1.
\]
\end{lemma}
\begin{proof}
对任意 \(1\le i,j\le k\)，有
\[
(B_kB_k^{\mathsf T})_{ij}
=\sum_{t=1}^{k}(B_k)_{it}(B_k)_{jt}
=\sum_{t=1}^{k}\mathbf{1}_{\{t\le i\}}\mathbf{1}_{\{t\le j\}}
=\sum_{t=1}^{\min(i,j)}1
=\min(i,j),
\]
故 \(K_k=B_kB_k^{\mathsf T}\)。
而 \(B_k\) 为单位下三角矩阵，故 \(\det(B_k)=1\)，从而 \(\det(K_k)=\det(B_k)^2=1\)。证毕。
\leanverified{minMatrix\_det\_at\_three}
\leanverified{minMatrix\_det\_at\_five}
\leanverified{minMatrix\_trace\_at\_small}
\leanverified{paper\_minMatrix\_small\_values\_package}
\end{proof}

\begin{corollary}[\texorpdfstring{\((4\sin^2)\)}{(4 sin^2)} 乘积恒等式与求和恒等式]\label{cor:pom-Kk-sine-product-sum}
令
\(
\phi_p:=\frac{(2p-1)\pi}{4k+2}
\)（\(p=1,\dots,k\)）。
则有严格恒等式
\[
\prod_{p=1}^{k}4\sin^2(\phi_p)=1,
\qquad
\sum_{p=1}^{k}\frac{1}{4\sin^2(\phi_p)}=\frac{k(k+1)}{2}.
\]
因此谱表示中的系数 \(4\) 受到行列式归一化刚性的约束：在保持整数 Gram 结构与 \(\det(K_k)=1\) 的前提下，
\((4\sin^2(\cdot))\) 不可再标度消去。
\leanverified{paper\_pom\_fold\_index\_seeds}
\end{corollary}
\begin{proof}
由引理 \ref{lem:pom-Kk-eigenvalues} 与注 \ref{rem:pom-Kk-factor4-spectral-rigidity}，\(K_k\) 的特征值为
\(
\lambda_p=(4\sin^2\phi_p)^{-1}
\)。
由引理 \ref{lem:pom-Kk-gram-det} 得 \(\det(K_k)=1=\prod_p\lambda_p\)，从而 \(\prod_p 4\sin^2(\phi_p)=1\)。
又 \(\mathrm{tr}(K_k)=\sum_{i=1}^k K_k(i,i)=\sum_{i=1}^k i=k(k+1)/2\)，而 \(\mathrm{tr}(K_k)=\sum_p\lambda_p\)，故得到求和恒等式。证毕。
\end{proof}

\begin{theorem}[Chebyshev 识别与 \texorpdfstring{\(\lambda/4\)}{lambda/4} 的整系数正规化]\label{thm:pom-Lk-chebyshev-charpoly}
记 \(L_k:=K_k^{-1}\in\ZZ^{k\times k}\)（其显式三对角形式见注 \ref{rem:pom-Kk-factor4-spectral-rigidity}）。
令 \(T_n\) 为第一类 Chebyshev 多项式（\(T_n(\cos\theta)=\cos(n\theta)\)），并定义
\[
P_k(y):=\frac{T_{2k+1}(\sqrt y)}{\sqrt y}\in\ZZ[y].
\]
则对任意 \(\lambda\in\CC\)，有特征多项式闭式
\[
\det(\lambda I-L_k)=(-1)^k\,P_k\!\Bigl(1-\frac{\lambda}{4}\Bigr).
\]
等价地，\(L_k\) 的谱点为
\[
\mu_p=4\sin^2\!\Bigl(\frac{(2p-1)\pi}{4k+2}\Bigr),\qquad p=1,\dots,k,
\]
并且线性变换 \(\lambda\mapsto 1-\lambda/4\) 是使 \(\det(\lambda I-L_k)\) 与 \(T_{2k+1}\) 发生整系数同构识别的自然正规化。
\leanverified{paper\_pom\_lk\_chebyshev\_charpoly}
\end{theorem}
\begin{proof}
由于 \(2k+1\) 为奇数，\(T_{2k+1}(z)\) 为奇多项式，故 \(T_{2k+1}(z)/z\) 为关于 \(z\) 的偶多项式，从而可唯一写成 \(P_k(z^2)\)；并且 \(T_n\) 具有整系数，故 \(P_k\in\ZZ[y]\)。
\par
由注 \ref{rem:pom-Kk-factor4-spectral-rigidity}，\(L_k\) 的谱满足
\(
\mu_p=4\sin^2(\phi_p)
\)
其中 \(\phi_p=\frac{(2p-1)\pi}{4k+2}\)。
对任意此类 \(\mu_p\)，令 \(y_p:=1-\mu_p/4=\cos^2(\phi_p)\)。
则 \(T_{2k+1}(\cos\phi_p)=\cos((2k+1)\phi_p)=\cos((2p-1)\pi/2)=0\)，从而
\(
P_k(y_p)=T_{2k+1}(\cos\phi_p)/\cos\phi_p=0
\)。
因此多项式 \(\lambda\mapsto P_k(1-\lambda/4)\) 的零点集恰为 \(\{\mu_p\}_{p=1}^k\)，与 \(\det(\lambda I-L_k)\) 的谱零点一致；
再由两者同为次数 \(k\) 的首一（除去符号）多项式，可得所示恒等式。证毕。
\end{proof}

\begin{corollary}[移位行列式闭式、自由能密度与重整化常数项]\label{cor:pom-Lk-shifted-det-free-energy}
\leanverified{paper\_pom\_lk\_shifted\_det\_free\_energy}
对任意实参数 \(t>-4\)，令 \(x:=\sqrt{1+t/4}>0\)。
则有严格恒等式
\[
\det(L_k+tI)=\det(I+tK_k)=P_k\!\Bigl(1+\frac{t}{4}\Bigr)=\frac{T_{2k+1}(x)}{x}.
\]
当 \(x>1\) 时，令 \(\eta:=\mathrm{arcosh}(x)\)，则
\[
\log\det(L_k+tI)=\log\cosh\bigl((2k+1)\eta\bigr)-\log x.
\]
并且存在自由能密度极限
\[
\lim_{k\to\infty}\frac{1}{2k+1}\log\det(L_k+tI)=\eta=\mathrm{arcosh}\!\Bigl(\sqrt{1+\frac{t}{4}}\Bigr).
\]
此外，定义重整化余项
\[
F_k(t):=\log\det(L_k+tI)-(2k+1)\,\mathrm{arcosh}\!\Bigl(\sqrt{1+\frac{t}{4}}\Bigr),
\]
则极限存在且为显式常数
\[
\lim_{k\to\infty}F_k(t)=-\log\!\Bigl(2\sqrt{1+\frac{t}{4}}\Bigr).
\]
\end{corollary}
\begin{proof}
由定理 \ref{thm:pom-Lk-chebyshev-charpoly}，将 \(\lambda=-t\) 代入并注意
\(
\det(-tI-L_k)=(-1)^k\det(L_k+tI)
\)
即可得到 \(\det(L_k+tI)=P_k(1+t/4)\)。
又由 \(K_kL_k=I\) 与引理 \ref{lem:pom-Kk-gram-det} 的 \(\det(K_k)=1\)，有
\(
I+tK_k=K_k(L_k+tI)
\)
从而 \(\det(I+tK_k)=\det(L_k+tI)\)。
当 \(x>1\) 时，\(T_{2k+1}(x)=\cosh((2k+1)\eta)\)，代回得到对数形式。
最后由 \(\log\cosh((2k+1)\eta)=(2k+1)\eta-\log 2+o(1)\) 得密度极限与常数项极限。证毕。
\end{proof}
% Determinant-polynomial rigidities for the mixed-boundary discrete Laplacian L_k.

\begin{proposition}[行列式多项式的系数闭式：\texorpdfstring{$\binom{k+j}{2j}$}{binom(k+j,2j)} 公式与正系数刚性]\label{prop:pom-Lk-det-coeff-binomial}
对 \(k\ge 0\) 定义
\[
D_k(t):=\det(L_k+tI),
\qquad
D_0(t):=1.
\]
则 \(D_k(t)\) 为次数 \(k\) 的整系数首一多项式，且具有完全显式的系数闭式
\[
\boxed{\;
D_k(t)=\sum_{j=0}^{k}\binom{k+j}{2j}\,t^{j}\; }.
\]
特别地，所有系数严格为正，并满足端点刚性
\[
[t^0]D_k=1,\qquad [t^{k-1}]D_k=2k-1,\qquad [t^{k}]D_k=1.
\]
\leanverified{detPoly\_eval\_zero}
\leanverified{detPoly\_monic}
\leanverified{detPoly\_natDegree}
\leanverified{detPoly\_coeff\_sub\_leading}
\leanverified{detPoly\_deriv\_eval\_zero\_double}
\leanverified{detPoly\_coeff\_binomial}
\leanverified{detPoly\_coeff\_pos}
\leanverified{paper\_pom\_lk\_det\_coeff\_binomial}
\end{proposition}

\begin{proof}
由推论 \ref{cor:pom-Lk-shifted-det-free-energy} 的双曲形式（取 \(t>0\)）
\[
D_k(t)=\frac{\cosh((2k+1)\eta(t))}{\cosh(\eta(t))},
\qquad
\eta(t)=\mathrm{arsinh}\sqrt{\frac{t}{4}},
\]
并用恒等式
\(
\cosh((n+2)\eta)=2\cosh(2\eta)\cosh(n\eta)-\cosh((n-2)\eta)
\)
（\(n\ge 1\)），得到递推
\[
D_{k+1}(t)=(t+2)D_k(t)-D_{k-1}(t),
\qquad
D_0(t)=1,\quad D_1(t)=t+1,
\]
其中使用了 \(2\cosh(2\eta(t))=t+2\)。
由于两边均为关于 \(t\) 的多项式，上述递推对一切 \(t\in\CC\) 成立。
\par
令
\(
F_k(t):=\sum_{j=0}^{k}\binom{k+j}{2j}\,t^j
\)。
直接计算得 \(F_0=1,F_1=t+1\)。
并且对任意 \(j\ge 0\) 有系数恒等式
\[
\binom{k+1+j}{2j}
\;=\;
2\binom{k+j}{2j}+\binom{k+j-1}{2j-2}-\binom{k+j-1}{2j},
\]
其证明只需两次应用 Pascal 公式：
\[
\binom{k+1+j}{2j}
=\binom{k+j}{2j}+\binom{k+j}{2j-1}
=\binom{k+j-1}{2j}+2\binom{k+j-1}{2j-1}+\binom{k+j-1}{2j-2}.
\]
因此 \(F_{k+1}=(t+2)F_k-F_{k-1}\)。
由同一递推与初值的唯一性，得 \(D_k\equiv F_k\)。
端点系数与严格正性由闭式直接给出。证毕。
\end{proof}

\begin{corollary}[Lucas 型闭式与二次域单位表示：\texorpdfstring{$t(4+t)$}{t(4+t)} 判别式的结构通道]\label{cor:pom-Lk-det-lucas-unit}
\leanverified{paper\_pom\_lk\_det\_lucas\_unit}
记判别式 \(\Delta(t):=t(4+t)\)，并在 \(\CC\setminus[-4,0]\) 上固定 \(\sqrt{\Delta(t)}\) 的解析分支（与 \(t>0\) 的正实分支相容）。
定义
\[
r(t):=\frac{t+2+\sqrt{\Delta(t)}}{2},
\qquad
r(t)^{-1}=\frac{t+2-\sqrt{\Delta(t)}}{2}.
\]
则 \(r(t)r(t)^{-1}=1\)、\(r(t)+r(t)^{-1}=t+2\)，并且对一切 \(k\ge 0\) 有闭式
\[
\boxed{\;
D_k(t)=\frac{r(t)^{\,k+1}+r(t)^{-k}}{r(t)+1}\; }.
\]
从而 \((D_k(t))_{k\ge 0}\) 满足二阶 Lucas 递推
\[
D_{k+1}(t)=(t+2)D_k(t)-D_{k-1}(t),
\qquad
D_0(t)=1,\quad D_1(t)=t+1.
\]
\leanverified{detPoly\_zero}
\leanverified{detPoly\_one}
\leanverified{detPoly\_succ\_succ}
\leanverified{detPoly\_two}
\leanverified{detPoly\_three}
\leanverified{detPoly\_eval\_two\_zero}
\leanverified{detPoly\_eval\_two\_one}
\leanverified{detPoly\_eval\_two\_recurrence}
\leanverified{detPoly\_eval\_neg\_one\_rec}
\leanverified{detPoly\_eval\_neg\_one\_periodic}
\leanverified{detPoly\_eval\_neg\_one\_zero}
\leanverified{detPoly\_eval\_neg\_one\_one}
\leanverified{detPoly\_eval\_neg\_one\_two}
\leanverified{detPoly\_eval\_neg\_one\_three}
\leanverified{detPoly\_eval\_neg\_one\_four}
\leanverified{detPoly\_eval\_neg\_one\_five}
\end{corollary}

\begin{proof}
由推论 \ref{cor:pom-Lk-area-law-qform}，对 \(t>0\) 有
\(
D_k(t)=q(t)^{-k}\frac{1+q(t)^{2k+1}}{1+q(t)}
\)
且 \(q(t)=e^{-2\eta(t)}\in(0,1)\)。
令 \(r(t):=q(t)^{-1}=e^{2\eta(t)}\)，则
\[
D_k(t)=\frac{r(t)^{\,k+1}+r(t)^{-k}}{r(t)+1}.
\]
另一方面，由 \(r+r^{-1}=e^{2\eta}+e^{-2\eta}=2\cosh(2\eta)=t+2\) 得
\(
r(t)=\frac{t+2+\sqrt{t(4+t)}}{2}
\)
（取与 \(t>0\) 相容的根号分支），故得到所示表述。
Lucas 递推由 \(r+r^{-1}=t+2\) 的一阶消元直接得到。证毕。
\end{proof}

\begin{proposition}[Cassini--Pell 守恒律：行列式 Wronskian 恒等式]\label{prop:pom-Lk-det-cassini-pell}
对一切 \(k\ge 1\) 与一切复参数 \(t\in\CC\) 成立严格恒等式
\[
\boxed{\;
D_{k+1}(t)D_{k-1}(t)-D_k(t)^2=t.\;}
\]
该恒等式等价于把 \((D_{k-1}(t),D_k(t))\) 嵌入二次曲线
\(
X^2-(t+2)XY+Y^2=t
\)
的整点轨道；在 \(t=1\) 时退化为奇指标 Fibonacci 行列式的 Cassini 恒等式。
\leanverified{detPoly\_cassini\_pell}
\leanverified{paper\_pom\_lk\_det\_cassini\_pell}
\leanverified{detPoly\_eval\_one}
\leanverified{fib\_odd\_cassini}
\leanverified{detPoly\_eval\_cassini\_gap}
\end{proposition}

\begin{proposition}[统一 Fibonacci Cassini 恒等式]\label{prop:pom-fib-cassini-unified}
对所有 \(n \in \mathbb{N}\)，在 \(\mathbb{Z}\) 中成立
\[
F_n \cdot F_{n+2} = F_{n+1}^2 + (-1)^{n+1}.
\]
即：当 \(n\) 为奇数时 \(F_n F_{n+2} = F_{n+1}^2 + 1\)；当 \(n\) 为偶数时 \(F_n F_{n+2} = F_{n+1}^2 - 1\)。
\leanverified{fib\_product\_cassini}
\end{proposition}

\begin{proof}
由推论 \ref{cor:pom-Lk-det-lucas-unit} 的递推，定义
\(
W_k:=D_{k+1}D_{k-1}-D_k^2
\)。
直接代入 \(D_{k+1}=(t+2)D_k-D_{k-1}\) 得
\[
W_k=\bigl((t+2)D_k-D_{k-1}\bigr)D_{k-1}-D_k^2
=D_k\bigl((t+2)D_{k-1}-D_{k-2}\bigr)-D_{k-1}^2
=W_{k-1},
\]
故 \(W_k\) 与 \(k\) 无关。
取 \(k=1\) 得
\(
W_1=D_2D_0-D_1^2=((t+2)(t+1)-1)- (t+1)^2=t
\)，从而对所有 \(k\ge 1\) 有 \(W_k=t\)。证毕。
\end{proof}

\begin{corollary}[严格对数凸性与比值单调：Cassini--Pell 不变量的曲率锁定]\label{cor:pom-Lk-det-logconvex-ratio}
当 \(t>0\) 时，由 \(D_k(t)>0\)（\(L_k+tI\succ 0\)）与命题 \ref{prop:pom-Lk-det-cassini-pell} 得对一切 \(k\ge 1\) 成立严格不等式
\[
D_k(t)^2< D_{k-1}(t)\,D_{k+1}(t),
\qquad
\frac{D_{k-1}(t)}{D_k(t)}>\frac{D_k(t)}{D_{k+1}(t)}.
\]
因此序列 \((D_k(t))_{k\ge 0}\) 严格对数凸，而比值序列 \(\bigl(D_{k-1}(t)/D_k(t)\bigr)_{k\ge 1}\) 严格单调递减；
在边界反射坐标 \(q_k(t):=D_{k-1}(t)/D_k(t)\) 下，这与 Schur--Riccati 递推的单调收敛一致（推论 \ref{cor:pom-Lk-riccati-error-by-determinants}）。
\leanverified{detPoly\_eval\_pos}
\leanverified{detPoly\_eval\_ratio\_gap\_pos}
\leanverified{detPoly\_eval\_strict\_log\_convex}
\leanverified{detPoly\_eval\_strict\_mono}
\end{corollary}
\begin{proof}
由命题 \ref{prop:pom-Lk-det-cassini-pell}，当 \(t>0\) 时有 \(D_{k+1}D_{k-1}-D_k^2=t>0\)，即得严格对数凸性。
再将该恒等式除以 \(D_kD_{k+1}>0\) 得
\[
\frac{D_{k-1}}{D_k}-\frac{D_k}{D_{k+1}}=\frac{t}{D_kD_{k+1}}>0,
\]
从而比值严格单调。证毕。
\end{proof}


\paragraph{确定性面积律：\texorpdfstring{\(\det(L_k+tI)\)}{det(L_k+tI)} 的 \(q\)-形式、Szeg\H{o} 分解与 resolvent/Green 闭式}
本段把推论 \ref{cor:pom-Lk-shifted-det-free-energy} 的行列式闭式进一步改写为 \(q\)-形式，并给出与符号积分完全对齐的 bulk--surface 分解，
以及 resolvent 迹与 Green 核矩阵元的双曲函数闭式。

\begin{corollary}[确定性“面积律”分解：\(q\)-形式与指数级有限尺寸余项]\label{cor:pom-Lk-area-law-qform}
\leanverified{paper\_pom\_lk\_area\_law\_qform}
取 \(t>0\)，并定义
\[
\eta(t):=\mathrm{arsinh}\sqrt{\frac{t}{4}}\qquad(\Re\eta\ge 0),
\qquad
q(t):=e^{-2\eta(t)}\in(0,1).
\]
则存在严格恒等式
\[
\det(L_k+tI)=\frac{\cosh\!\bigl((2k+1)\eta(t)\bigr)}{\cosh(\eta(t))}
=q(t)^{-k}\,\frac{1+q(t)^{2k+1}}{1+q(t)}.
\]
从而
\[
\log\det(L_k+tI)=2k\,\eta(t)-\log(1+q(t))+\log\bigl(1+q(t)^{2k+1}\bigr),
\]
其中最后一项为完全显式的有限尺寸余项，并满足
\[
0<\log\bigl(1+q(t)^{2k+1}\bigr)<q(t)^{2k+1}.
\]
\end{corollary}
\begin{proof}
由推论 \ref{cor:pom-Lk-shifted-det-free-energy}，令 \(x=\sqrt{1+t/4}=\cosh(\eta(t))\) 则
\[
\det(L_k+tI)=\frac{T_{2k+1}(x)}{x}=\frac{\cosh((2k+1)\eta(t))}{\cosh(\eta(t))}.
\]
再令 \(q=e^{-2\eta}\)，把 \(\cosh((2k+1)\eta)/\cosh(\eta)\) 写成指数形式即可得到
\(
q^{-k}(1+q^{2k+1})/(1+q)
\)；
取对数即得所示分解。
最后，由 \(q\in(0,1)\) 知 \(\log(1+q^{2k+1})\in(0,q^{2k+1})\)。证毕。
\end{proof}

\begin{proposition}[强 Szeg\H{o} 型分解的可积闭式：符号 \(a_t(\theta)=1+\frac{t}{4\sin^2(\theta/2)}\)]\label{prop:pom-Kk-strong-szego-bulk-surface}
\leanverified{paper\_pom\_kk\_strong\_szego\_bulk\_surface}
由 \(\det(K_k)=1\)（引理 \ref{lem:pom-Kk-gram-det}）知 \(\det(I+tK_k)=\det(L_k+tI)\)。
对 \(t>0\) 定义符号函数
\[
a_t(\theta):=1+\frac{t}{4\sin^2(\theta/2)}\qquad(\theta\in(0,\pi)).
\]
则存在严格分解
\[
\log\det(I+tK_k)
=k\cdot \frac{1}{\pi}\int_{0}^{\pi}\log a_t(\theta)\,\dd\theta
\;+\;\log\frac{1+q(t)^{2k+1}}{1+q(t)},
\]
并且体相积分项可显式化为
\[
\frac{1}{\pi}\int_{0}^{\pi}\log\!\Bigl(1+\frac{t}{4\sin^2(\theta/2)}\Bigr)\,\dd\theta
=2\,\eta(t).
\]
因此 \(\log\det(I+tK_k)\) 的 bulk 密度恰为 \(2\eta(t)\)，而符号中分母的系数 \(4\) 由该闭式同构与 Chebyshev--双曲参数 \(\eta\) 的对齐结构性强制固定。
\end{proposition}
\begin{proof}
由推论 \ref{cor:pom-Lk-area-law-qform}，
\(
\log\det(I+tK_k)=2k\,\eta(t)+\log\frac{1+q^{2k+1}}{1+q}
\)。
故只需证明
\(
\frac{1}{\pi}\int_0^\pi \log a_t(\theta)\,\dd\theta=2\eta(t)
\)。
设
\(
J(t):=\frac{1}{\pi}\int_0^\pi \log\bigl(1+\frac{t}{4\sin^2(\theta/2)}\bigr)\,\dd\theta
\)。
对 \(t>0\) 可在积分号下求导得
\[
J'(t)=\frac{1}{\pi}\int_0^\pi \frac{1}{t+4\sin^2(\theta/2)}\,\dd\theta
=\frac{1}{\pi}\int_0^\pi \frac{1}{t+2(1-\cos\theta)}\,\dd\theta
=\frac{1}{\sqrt{t(4+t)}},
\]
其中最后一步为标准积分恒等式（亦可由复积分或对 \(t\) 的 Green 函数解析解得到）。
另一方面，\(\eta(t)=\mathrm{arsinh}\sqrt{t/4}\) 满足 \(\eta'(t)=1/(2\sqrt{t(4+t)})\)，故 \((2\eta)'=J'\)。
又 \(J(0)=0=2\eta(0)\)，从而 \(J(t)=2\eta(t)\)。证毕。
\end{proof}

\begin{corollary}[解析可审计的 resolvent 迹：\texorpdfstring{\(\tr(L_k+tI)^{-1}\)}{tr((L_k+tI)^{-1})} 的双曲闭式]\label{cor:pom-Lk-resolvent-trace-hyperbolic}
取 \(t>0\) 并沿用 \(\eta(t)\) 记号。
则 resolvent 迹满足严格恒等式
\[
\tr(L_k+tI)^{-1}
=\frac{\partial}{\partial t}\log\det(L_k+tI)
=\frac{2k+1}{2\sqrt{t(4+t)}}\,\tanh\!\bigl((2k+1)\eta(t)\bigr)-\frac{1}{2(4+t)}.
\]
特别地，迹密度极限存在且为
\[
\lim_{k\to\infty}\frac{1}{2k+1}\tr(L_k+tI)^{-1}
=\frac{1}{2\sqrt{t(4+t)}}.
\]
\end{corollary}
\begin{proof}
由推论 \ref{cor:pom-Lk-shifted-det-free-energy}，
\(
\log\det(L_k+tI)=\log\cosh((2k+1)\eta)-\log\cosh\eta
\)（其中 \(\eta=\eta(t)\)）。
对 \(t\) 求导并用 \(\eta'(t)=1/(2\sqrt{t(4+t)})\) 得
\[
\frac{\partial}{\partial t}\log\det(L_k+tI)
=\eta'(t)\Bigl((2k+1)\tanh((2k+1)\eta)-\tanh\eta\Bigr)
=\frac{2k+1}{2\sqrt{t(4+t)}}\tanh((2k+1)\eta)-\frac{1}{2(4+t)},
\]
其中最后一步用到 \(\tanh(\eta)=\sqrt{t/(4+t)}\)。
而 \(\tr(L_k+tI)^{-1}=\partial_t\log\det(L_k+tI)\) 为一般恒等式。
取 \(k\to\infty\) 且用 \(\eta>0\Rightarrow\tanh((2k+1)\eta)\to 1\) 即得极限。证毕。
\end{proof}

\begin{proposition}[Green 核矩阵元的完全闭式：\texorpdfstring{\(((L_k+tI)^{-1}_{ij})\)}{((L_k+tI)^{-1}_{ij})} 的双曲乘积公式]\label{prop:pom-Lk-green-kernel-entries}
\leanverified{paper\_pom\_Lk\_green\_kernel\_entries}
取 \(t>0\) 并记 \(\eta=\eta(t)\)。
则对任意 \(1\le i\le j\le k\)，逆矩阵矩阵元满足严格闭式
\[
(L_k+tI)^{-1}_{ij}
=\frac{\sinh(2i\eta)\,\cosh\!\bigl((2(k-j)+1)\eta\bigr)}
{\sinh(2\eta)\,\cosh\!\bigl((2k+1)\eta\bigr)},
\qquad
(L_k+tI)^{-1}_{ji}=(L_k+tI)^{-1}_{ij}.
\]
特别地，
\[
(L_k+tI)^{-1}_{11}=\frac{\cosh((2k-1)\eta)}{\cosh((2k+1)\eta)},
\qquad
(L_k+tI)^{-1}_{1k}=\frac{\cosh(\eta)}{\cosh((2k+1)\eta)}.
\]
\end{proposition}
\begin{proof}
令 \(A:=L_k+tI\)。它是对称三对角矩阵，除末端外对角元恒为 \(2+t\)，末端对角元为 \(1+t\)，非对角元恒为 \(-1\)。
对 \(r\ge 0\) 定义 leading 主子式
\(
D_r:=\det(A_{[1..r,\,1..r]})
\)
并约定 \(D_0:=1\)。则对 \(1\le r\le k-1\) 有递推 \(D_r=(2+t)D_{r-1}-D_{r-2}\)，从而
\[
D_{i-1}=\frac{\sinh(2i\eta)}{\sinh(2\eta)}.
\]
同理，定义 trailing 主子式
\(
E_r:=\det(A_{[k-r+1..k,\,k-r+1..k]})
\)
并约定 \(E_0:=1\)、\(E_1:=1+t\)。递推与端点条件解出
\[
E_{k-j}=\frac{\cosh\!\bigl((2(k-j)+1)\eta\bigr)}{\cosh(\eta)}.
\]
对三对角矩阵的 Cramer--Gantmacher 公式（亦即一维差分 Green 函数的标准表示）给出：当 \(i\le j\) 时
\[
A^{-1}_{ij}=\frac{D_{i-1}\,E_{k-j}}{D_k},
\]
其中 \(D_k=\det(A)=\cosh((2k+1)\eta)/\cosh(\eta)\) 由推论 \ref{cor:pom-Lk-shifted-det-free-energy} 给出。
代入即得所示闭式。证毕。
\end{proof}

\begin{corollary}[指数聚类与相关长度：质量参数 \(t\) 的严格“间隙律”]\label{cor:pom-Lk-exponential-clustering-gaplaw}
\leanverified{paper\_pom\_lk\_exponential\_clustering\_gaplaw}
取 \(t>0\) 并令 \(\eta=\eta(t)>0\)。则对任意 \(1\le i\le j\le k\) 有一致上界
\[
(L_k+tI)^{-1}_{ij}\ \le\ \frac{1}{\sinh(2\eta)}\,e^{-2\eta(j-i)}.
\]
从而可定义相关长度
\[
\xi(t):=\frac{1}{2\eta(t)}=\frac{1}{2\,\mathrm{arsinh}\sqrt{t/4}},
\]
并得到严格指数聚类：在固定 \(t>0\) 下，Green 关联随格距 \(j-i\) 以速率 \(e^{-(j-i)/\xi(t)}\) 衰减。
\end{corollary}
\begin{proof}
由命题 \ref{prop:pom-Lk-green-kernel-entries}，
\[
(L_k+tI)^{-1}_{ij}
=\frac{\sinh(2i\eta)\,\cosh((2(k-j)+1)\eta)}{\sinh(2\eta)\,\cosh((2k+1)\eta)}.
\]
用 \(\sinh u\le \tfrac12 e^{u}\) 与 \(\cosh u\le e^{u}\)，以及 \(\cosh((2k+1)\eta)\ge \tfrac12 e^{(2k+1)\eta}\)，得到
\[
(L_k+tI)^{-1}_{ij}
\le
\frac{(\tfrac12 e^{2i\eta})\cdot (e^{(2(k-j)+1)\eta})}{\sinh(2\eta)\cdot (\tfrac12 e^{(2k+1)\eta})}
=\frac{1}{\sinh(2\eta)}\,e^{-2\eta(j-i)}.
\]
相关长度表述由 \(\xi=1/(2\eta)\) 的重写立得。证毕。
\end{proof}

\begin{corollary}[半无限极限的边界固定点：反射系数 \(q(t)\) 的代数闭式与 Riccati 方程]\label{cor:pom-Lk-boundary-fixed-point-riccati}
\leanverified{paper\_pom\_lk\_boundary\_fixed\_point\_riccati}
取 \(t>0\) 并令 \(q(t)=e^{-2\eta(t)}\)。
则边界 Green 值存在且满足
\[
\lim_{k\to\infty}(L_k+tI)^{-1}_{11}=q(t).
\]
并且 \(q(t)\) 具有闭式
\[
q(t)=\Bigl(\frac{\sqrt{4+t}-\sqrt{t}}{2}\Bigr)^{2}
=1+\frac{t}{2}-\frac{1}{2}\sqrt{t(4+t)},
\]
且是二次方程的唯一 \((0,1)\) 解：
\[
q^2-(2+t)q+1=0.
\]
\end{corollary}
\begin{proof}
由命题 \ref{prop:pom-Lk-green-kernel-entries}，\((L_k+tI)^{-1}_{11}=\cosh((2k-1)\eta)/\cosh((2k+1)\eta)\)，令 \(k\to\infty\) 得极限为 \(e^{-2\eta}=q\)。
又 \(\eta=\mathrm{arsinh}(\sqrt{t}/2)\)，故 \(e^{-\eta}=\sqrt{1+t/4}-\sqrt{t}/2\)，从而
\[
q=e^{-2\eta}=\Bigl(\sqrt{1+\frac{t}{4}}-\frac{\sqrt{t}}{2}\Bigr)^2
=\Bigl(\frac{\sqrt{4+t}-\sqrt{t}}{2}\Bigr)^2,
\]
展开即得第二式。
最后，由 \(2\cosh(2\eta)=e^{2\eta}+e^{-2\eta}=q^{-1}+q=2+t\) 得 \(q^2-(2+t)q+1=0\)，并由 \(q\in(0,1)\) 的分支选取唯一性得结论。证毕。
\end{proof}

\begin{corollary}[表面自由能与边界响应：常数项闭式与可观测导数]\label{cor:pom-Lk-surface-free-energy}
\leanverified{paper\_pom\_lk\_surface\_free\_energy}
取 \(t>0\) 并记 \(\eta=\eta(t)\)、\(q=q(t)\)。
定义 bulk 自由能密度与表面自由能
\[
f(t):=\lim_{k\to\infty}\frac{1}{k}\log\det(L_k+tI)=2\eta(t),
\qquad
g(t):=\lim_{k\to\infty}\Bigl(\log\det(L_k+tI)-k f(t)\Bigr)=-\log(1+q(t)).
\]
则对任意 \(k\ge 1\) 有严格恒等式
\[
\log\det(L_k+tI)=k f(t)+g(t)+\log\bigl(1+q(t)^{2k+1}\bigr),
\]
且表面项的导数可由边界响应给出：
\[
g'(t)=\lim_{k\to\infty}\Bigl(\tr(L_k+tI)^{-1}-k f'(t)\Bigr)
=\frac{q(t)}{1+q(t)}\cdot \frac{1}{\sqrt{t(4+t)}}.
\]
\end{corollary}
\begin{proof}
由推论 \ref{cor:pom-Lk-area-law-qform} 的精确分解
\(
\log\det=2k\eta-\log(1+q)+\log(1+q^{2k+1})
\)，取极限即得 \(f(t)=2\eta\) 与 \(g(t)=-\log(1+q)\) 及所示恒等式。
对导数部分，注意 \(f'(t)=2\eta'(t)=1/\sqrt{t(4+t)}\)，并由推论 \ref{cor:pom-Lk-resolvent-trace-hyperbolic} 得
\[
\tr(L_k+tI)^{-1}-k f'(t)
=\frac{2k+1}{2\sqrt{t(4+t)}}\tanh((2k+1)\eta)-\frac{1}{2(4+t)}-\frac{k}{\sqrt{t(4+t)}}.
\]
令 \(k\to\infty\)（\(\eta>0\Rightarrow\tanh((2k+1)\eta)\to 1\)）得到极限
\(
\frac{1}{2\sqrt{t(4+t)}}-\frac{1}{2(4+t)}
\)。
而由推论 \ref{cor:pom-Lk-boundary-fixed-point-riccati} 的 \(q\)-闭式可化简为
\(
\frac{q}{1+q}\cdot \frac{1}{\sqrt{t(4+t)}}
\)，与 \(g'(t)=-(q'/(1+q))\)（\(q'=-(2\eta')q\)）一致。证毕。
\end{proof}

\begin{corollary}[表面自由能的严格凹性：二阶导数闭式与符号刚性]\label{cor:pom-Lk-surface-free-energy-strict-concavity}
\leanverified{paper\_pom\_lk\_surface\_free\_energy\_strict\_concavity}
沿用推论 \ref{cor:pom-Lk-surface-free-energy} 的记号。
则对任意 \(t>0\) 有二阶导数闭式
\[
\boxed{\;
g''(t)
=-\frac{q(t)}{(1+q(t))^2\,t(4+t)}
-\frac{q(t)(t+2)}{(1+q(t))\,[t(4+t)]^{3/2}}\;<\;0\;}
\]
其中 \(q(t)=1+\frac{t}{2}-\frac12\sqrt{t(4+t)}\in(0,1)\) 为边界固定点（推论 \ref{cor:pom-Lk-boundary-fixed-point-riccati}）。
因此 \(g\) 在整个正轴上严格凹，其凹性由 \(q(t)\) 的代数闭式强制。
\end{corollary}
\begin{proof}
由 \(g(t)=-\log(1+q(t))\) 与推论 \ref{cor:pom-Lk-surface-free-energy} 的导数闭式
\(
g'(t)=\frac{q(t)}{1+q(t)}\cdot\frac{1}{\sqrt{t(4+t)}}
\)
直接求导，并用 \(q(t)=1+\frac{t}{2}-\frac12\sqrt{t(4+t)}\) 化简，即得所示表达式。
由于 \(q(t)\in(0,1)\) 且 \(t(4+t)>0\)，右端两项均严格为负，故 \(g''(t)<0\)。证毕。
\end{proof}

\begin{remark}[有限窗审计：确定性面积律分解、resolvent 迹与 Green 核矩阵元]\label{rem:pom-fence-green-kernel-area-law-audit}
脚本在有限窗上核验：推论 \ref{cor:pom-Lk-area-law-qform} 的 \(q\)-形式分解与指数小余项界、
命题 \ref{prop:pom-Kk-strong-szego-bulk-surface} 的符号积分闭式、
推论 \ref{cor:pom-Lk-resolvent-trace-hyperbolic} 的 resolvent 迹闭式、
命题 \ref{prop:pom-Lk-green-kernel-entries} 的矩阵元闭式与推论 \ref{cor:pom-Lk-exponential-clustering-gaplaw} 的指数聚类上界，
并对边界固定点 \(q(t)\) 与表面自由能导数 \(g'(t)\) 给出数值一致性证书。
审计表如下：
\begin{table}[H]
\centering
\scriptsize
\setlength{\tabcolsep}{6pt}
\caption{确定性面积律/Green-resolvent 闭式的有限窗审计（取 $t=1$）：核验 $q$-形式行列式分解（推论~\ref{cor:pom-Lk-area-law-qform}）、符号积分闭式（命题~\ref{prop:pom-Kk-strong-szego-bulk-surface}）、resolvent 迹闭式（推论~\ref{cor:pom-Lk-resolvent-trace-hyperbolic}）、Green 矩阵元闭式与指数聚类上界（命题~\ref{prop:pom-Lk-green-kernel-entries}、推论~\ref{cor:pom-Lk-exponential-clustering-gaplaw}），以及边界固定点 $q(t)$ 与表面项导数 $g'(t)$ 的一致性（推论~\ref{cor:pom-Lk-boundary-fixed-point-riccati}、\ref{cor:pom-Lk-surface-free-energy}）。}
\label{tab:fence_green_kernel_area_law_resolvent_audit}
\begin{tabular}{r r r r r r r r}
\toprule
$k$ & $|\Delta\log\det|$ & $|\Delta\,\mathrm{bulk}|$ & $|\Delta\,\mathrm{tr}|$ & $\max|\Delta G_{ij}|$ & $\max\rho$ & $|G_{11}-q|$ & $|\Delta g'|$\\
\midrule
10 & 0.00e+00 & 1.20e-04 & 8.88e-16 & 2.78e-16 & 0.590170 & 3.33e-16 & 3.91e-15\\
20 & 3.55e-15 & 1.20e-04 & 0.00e+00 & 3.89e-16 & 0.676275 & 3.33e-16 & 3.91e-15\\
40 & 7.11e-15 & 1.20e-04 & 3.55e-15 & 5.55e-16 & 0.500000 & 3.33e-16 & 3.91e-15\\
200 & 0.00e+00 & 1.20e-04 & 0.00e+00 & -- & 0.500000 & 3.33e-16 & 7.47e-15\\
\bottomrule
\end{tabular}
\end{table}

\end{remark}


\paragraph{黄金耦合与 Fisher 零点：\texorpdfstring{\(t=1\)}{t=1} 的 Fibonacci 算术化、零点反正弦律与 Riccati 递推}
本段把质量参数 \(t\) 的“黄金刻度”与离散行列式/Green 核的可审计闭式联结起来：
一方面在 \(t=1\) 处得到完全算术化的 Fibonacci 行列式与两点关联；
另一方面把 \(\det(L_k+tI)\) 的零点（Fisher 零点）与 \([-4,0]\) 上的反正弦律严格对齐，并给出零点的均匀化与边界 Riccati 递推的精确收敛率。

\begin{corollary}[“黄金耦合”唯一性：\texorpdfstring{\(t=1\)}{t=1} 的等价刻画]\label{cor:pom-Lk-golden-coupling-unique}
设 \(\varphi=\frac{1+\sqrt5}{2}\)，并对 \(t>0\) 定义
\[
\eta(t):=\mathrm{arsinh}\sqrt{\frac{t}{4}}\qquad(\Re\eta\ge 0).
\]
则下述命题等价，并给出同一唯一正解 \(t_\star=1\)：
\[
\eta(t_\star)=\log\varphi,
\qquad
q(t_\star):=e^{-2\eta(t_\star)}=\varphi^{-2},
\qquad
\lim_{k\to\infty}\frac{1}{k}\log\det(L_k+t_\star I)=2\log\varphi.
\]
因此 \(t=1\) 是使谱自由能密度严格匹配黄金移位熵刻度 \(2\log\varphi\) 的唯一正参数点。
\leanverified{sinh\_log\_phi\_eq\_half}
\leanverified{golden\_coupling\_q\_value}
\leanverified{paper\_pom\_golden\_coupling\_uniqueness}
\end{corollary}
\begin{proof}
由定义，\(\eta(t)\) 在 \(t>0\) 上严格递增，因此方程 \(\eta(t)=\log\varphi\) 至多有一解。
又
\(
\eta(t)=\mathrm{arsinh}(\sqrt{t}/2)
\)
等价于 \(\sqrt{t}/2=\sinh(\eta)\)。
取 \(\eta=\log\varphi\)，利用
\(
\sinh(\log\varphi)=\frac{\varphi-\varphi^{-1}}{2}=\frac12
\)
得到 \(t=1\)。
第二个等价式由 \(q(t)=e^{-2\eta(t)}\) 的定义立得。
第三个等价式由推论 \ref{cor:pom-Lk-surface-free-energy} 的 \(f(t)=\lim_{k\to\infty}\frac1k\log\det(L_k+tI)=2\eta(t)\) 给出。证毕。
\end{proof}

\begin{corollary}[黄金耦合处的 Fibonacci 行列式与 Green 核全显式：分母 \texorpdfstring{\(F_{2k+1}\)}{F_{2k+1}} 的统一化]\label{cor:pom-Lk-t1-fibonacci-det-green}
仍以 \(F_0=0,F_1=1\) 为 Fibonacci 数列。令 \(A_k:=L_k+I\in\ZZ^{k\times k}\)。
则
\[
\det(A_k)=F_{2k+1}.
\]
并且其逆矩阵矩阵元对 \(1\le i\le j\le k\) 满足精确闭式
\[
(A_k^{-1})_{ij}
=\frac{F_{2i}\,F_{2(k-j)+1}}{F_{2k+1}},
\qquad
(A_k^{-1})_{ji}=(A_k^{-1})_{ij}.
\]
等价地，伴随矩阵在对象层发生分离变量因子化：
\[
\operatorname{adj}(A_k)_{ij}=F_{2\min(i,j)}\,F_{2(k-\max(i,j))+1}.
\]
特别地，
\[
(A_k^{-1})_{11}=\frac{F_{2k-1}}{F_{2k+1}},
\qquad
(A_k^{-1})_{1k}=\frac{1}{F_{2k+1}}.
\]
\leanverified{fenceDet\_eq\_detPoly\_eval\_one}
\leanverified{fenceDet\_partial\_sum}
\leanverified{fenceDet\_partial\_sum\_5}
\leanverified{fenceDet\_partial\_sum\_pos}
\leanverified{paper\_fenceDet\_partial\_sum\_package}
\leanverified{fenceDet\_succ\_sub}
\leanverified{fenceDet\_succ\_eq\_add}
\leanverified{paper\_fenceDet\_diff\_package}
\leanverified{fenceDet\_mul\_succ}
\leanverified{paper\_fenceDet\_product\_package}
\leanverified{fenceDet\_eight}
\leanverified{fenceDet\_nine}
\leanverified{fenceDet\_ten}
\leanverified{prime\_1597}
\leanverified{paper\_fenceDet\_values\_extended}
\leanverified{paper\_pom\_Lk\_t1\_fibonacci\_det\_green}
\end{corollary}
\begin{proof}
由推论 \ref{cor:pom-Lk-shifted-det-free-energy}，对 \(t=1\) 有
\[
\det(A_k)=\det(L_k+I)=\frac{\cosh((2k+1)\eta(1))}{\cosh(\eta(1))}.
\]
而 \(\eta(1)=\mathrm{arsinh}(1/2)=\log\varphi\)，并注意 \(\cosh(\log\varphi)=\frac{\varphi+\varphi^{-1}}{2}=\frac{\sqrt5}{2}\)。
因此
\[
\det(A_k)=\frac{\cosh((2k+1)\log\varphi)}{\cosh(\log\varphi)}
=\frac{(\varphi^{2k+1}+\varphi^{-(2k+1)})/2}{\sqrt5/2}
=\frac{\varphi^{2k+1}+\varphi^{-(2k+1)}}{\sqrt5}
=F_{2k+1}.
\]
矩阵元闭式由命题 \ref{prop:pom-Lk-green-kernel-entries} 取 \(t=1\) 得到：
\[
(A_k^{-1})_{ij}
=\frac{\sinh(2i\log\varphi)\,\cosh((2(k-j)+1)\log\varphi)}{\sinh(2\log\varphi)\,\cosh((2k+1)\log\varphi)}.
\]
利用恒等式
\(
\sinh(2n\log\varphi)=\frac{\sqrt5}{2}F_{2n}
\)
与
\(
\cosh((2m+1)\log\varphi)=\frac{\sqrt5}{2}F_{2m+1}
\)
（由 Binet 公式直接化简）即得所示 Fibonacci 比。
伴随矩阵公式由 \(\operatorname{adj}(A_k)=\det(A_k)\,A_k^{-1}\) 立得。证毕。
\leanverified{paper\_fenceDet\_values\_and\_strict\_mono}
\end{proof}

\begin{remark}[“面积项 + 常数项 + 指数小修正”在 Fibonacci 上的精确实现]\label{rem:pom-Lk-fibonacci-area-law-exact}
在黄金耦合 \(t=1\) 下，由 Binet 公式对奇指标的精确写法
\[
F_{2k+1}=\frac{\varphi^{2k+1}}{\sqrt5}\Bigl(1+\varphi^{-4k-2}\Bigr)
\]
得到不经渐近的精确分解
\[
\log F_{2k+1}=(2k+1)\log\varphi-\frac12\log 5+\log\bigl(1+\varphi^{-4k-2}\bigr),
\qquad
0<\log\bigl(1+\varphi^{-4k-2}\bigr)<\varphi^{-4k-2}.
\]
因此，“面积律主项 + 普适常数项 + 指数小有限尺寸修正”在本模型中可视为 Fibonacci 的严格代数恒等式，而非经验级近似。
\end{remark}

\begin{corollary}[边界熵导数与 Parry 质量的精确耦合：\texorpdfstring{\(t=1\)}{t=1} 处的同源闭合]\label{cor:pom-Lk-t1-parry-surface-derivative}
令黄金均值 shift 的 Parry 最大熵测度在符号 \(1\) 上的质量为
\(
\pi(1)=\frac{1}{\varphi^2+1}
\)
（注 \ref{prop:parry-measure-golden}）。
在谱侧，令 \(q(t)=e^{-2\eta(t)}\) 并取表面自由能 \(g(t)=-\log(1+q(t))\)（推论 \ref{cor:pom-Lk-surface-free-energy}）。
则在黄金耦合点 \(t=1\) 有严格恒等式
\[
\frac{q(1)}{1+q(1)}=\frac{\varphi^{-2}}{1+\varphi^{-2}}=\frac{1}{\varphi^2+1}=\pi(1),
\]
并且
\[
g'(1)=\pi(1)\cdot\frac{1}{\sqrt{1\cdot(4+1)}}=\frac{\pi(1)}{\sqrt5}.
\]
\leanverified{parry\_mass\_identity}
\leanverified{parry\_surface\_coupling}
\leanverified{t\_times\_4\_plus\_t\_at\_one}
\leanverified{paper\_pom\_parry\_surface\_derivative\_seeds}
\leanverified{paper\_pom\_parry\_surface\_derivative\_package}
\end{corollary}
\begin{proof}
由推论 \ref{cor:pom-Lk-golden-coupling-unique} 知 \(q(1)=\varphi^{-2}\)，代入即得第一式。
第二式由推论 \ref{cor:pom-Lk-surface-free-energy} 的
\(
g'(t)=\frac{q(t)}{1+q(t)}\cdot \frac{1}{\sqrt{t(4+t)}}
\)
在 \(t=1\) 处取值立得。证毕。
\leanverified{paper\_pom\_Lk\_t1\_parry\_surface\_derivative}
\end{proof}

\begin{theorem}[Fisher 零点的极限分布：\texorpdfstring{\(\det(L_k+tI)\)}{det(L_k+tI)} 的零点测度收敛到 \texorpdfstring{\([-4,0]\)}{[-4,0]} 上的反正弦律]\label{thm:pom-Lk-fisher-zeros-arcsine}
把 \(P_k(t):=\det(L_k+tI)\) 视为关于 \(t\) 的次数 \(k\) 多项式，其零点为
\[
t_p=-4\sin^2\!\Bigl(\frac{(2p-1)\pi}{4k+2}\Bigr)\in(-4,0),
\qquad p=1,\dots,k.
\]
令零点经验测度
\(
\omega_k:=\frac1k\sum_{p=1}^k\delta_{t_p}
\)。
则 \(\omega_k\) 在弱拓扑下收敛到 \([-4,0]\) 上的反正弦分布 \(\omega\)，其密度为
\[
\dd\omega(t)=\frac{1}{\pi\sqrt{(-t)(4+t)}}\,\dd t,\qquad t\in(-4,0).
\]
并且对任意 \(t\in\CC\setminus[-4,0]\)，归一化对数势存在且为闭式
\[
\lim_{k\to\infty}\frac{1}{k}\log\bigl|\det(L_k+tI)\bigr|
=2\,\Re\,\mathrm{arsinh}\sqrt{\frac{t}{4}}.
\]
\leanverified{paper\_pom\_fisher\_zeros\_arcsine\_seeds}
\leanverified{paper\_pom\_fisher\_zeros\_arcsine\_package}
\leanverified{paper\_pom\_fisher\_zeros\_arcsine}
\end{theorem}
\begin{proof}
由于 \(P_k(t)=\prod_{p=1}^k(\mu_p+t)\) 且 \(\mu_p\) 为 \(L_k\) 的特征值（定理 \ref{thm:pom-Lk-arcsine-law}），零点满足 \(t_p=-\mu_p\)。
故 \(\omega_k\) 是经验谱测度 \(\nu_k\) 在映射 \(t=-\mu\) 下的推前；由定理 \ref{thm:pom-Lk-arcsine-law} 立得极限密度为所示反正弦律。

对对数势部分，由推论 \ref{cor:pom-Lk-shifted-det-free-energy} 的 Chebyshev 形式
\(
\det(L_k+tI)=T_{2k+1}(\sqrt{1+t/4})/\sqrt{1+t/4}
\)
并用 \(T_n(z)=\frac12(w^n+w^{-n})\)（其中 \(w=z+\sqrt{z^2-1}\)，在 \(\CC\setminus[-1,1]\) 上解析）可得
\[
\frac{1}{k}\log|\det(L_k+tI)|
=\frac{2k+1}{k}\Re\log w(t)+o(1),
\]
其中 \(w(t)=\sqrt{1+t/4}+\sqrt{t/4}=\exp(\mathrm{arsinh}\sqrt{t/4})\)（选取与 \(t>0\) 相容的解析延拓）。
令 \(k\to\infty\) 即得极限 \(2\Re\,\mathrm{arsinh}\sqrt{t/4}\)。证毕。
\end{proof}

\begin{remark}[零点的共形均匀化：Joukowsky 映射下的“奇根统一”]\label{rem:pom-Lk-fisher-zeros-uniformization}
引入均匀化参数 \(w\in\CC^\times\) 使
\[
w+w^{-1}=t+2
\qquad\Longleftrightarrow\qquad
t=w+w^{-1}-2.
\]
则 \(P_k(t)=0\) 等价于
\[
w^{2k+1}=-1.
\]
因此零点集合是“\(-1\) 的奇阶单位根”在 Joukowsky 映射 \(t=w+w^{-1}-2\) 下的像；
当 \(|w|=1\) 时，像域恰为线段 \([-4,0]\)，并把单位圆上的均匀离散点推前为定理 \ref{thm:pom-Lk-fisher-zeros-arcsine} 的反正弦律极限。
该均匀化同时给出端点尺度的几何来源：\([-4,0]\) 是单位圆在 \(w+w^{-1}-2\) 下的精确像域端点。
\end{remark}

\begin{proposition}[临界端点的平方根型分支：\texorpdfstring{$t=0$}{t=0} 与 \texorpdfstring{$t=-4$}{t=-4} 的 Puiseux 指数刚性]\label{prop:pom-Lk-branch-puiseux}
\leanverified{paper\_pom\_lk\_branch\_puiseux}
把
\[
f(t):=\lim_{k\to\infty}\frac{1}{k}\log\det(L_k+tI)
=2\,\mathrm{arsinh}\sqrt{\frac{t}{4}}
\]
视为 \(t\in\CC\setminus[-4,0]\) 上的解析函数（固定与 \(t>0\) 相容的 \(\sqrt{t(4+t)}\) 分支）。
则 \(t=0\) 与 \(t=-4\) 是仅有的平方根型分支点（对应判别式 \(\Delta(t)=t(4+t)\) 的两零点）。
并且在外域极限 \(t\to 0\) 下有 Puiseux 展开
\[
\boxed{\;
f(t)=\sqrt{t}-\frac{t^{3/2}}{24}+O(t^{5/2})\;}
\]
（分支由 \(\sqrt{t}\) 的选取确定）。
另一方面，令表面自由能 \(g(t):=-\log(1+e^{-2\eta(t)})\)（其中 \(\eta(t)=\mathrm{arsinh}\sqrt{t/4}\)，见推论 \ref{cor:pom-Lk-surface-free-energy}），则同一极限下
\[
\boxed{\;
g(t)=-\log 2+\frac{\sqrt{t}}{2}+O(t)\;}
\]
从而 bulk 与表面项的非解析性均由同一 Puiseux 指数 \(1/2\) 控制；并且第二分支点 \(t=-4\) 的位置恰为 \(\Delta(t)\) 的另一零点，等价于分支点间距的规范化常数 \(4\)。
\end{proposition}

\begin{proof}
令 \(z=\sqrt{t}/2\)，则 \(\eta(t)=\mathrm{arsinh}(z)\)。
由级数 \(\mathrm{arsinh}(z)=z-\frac{z^3}{6}+O(z^5)\)（\(z\to 0\)）得
\[
f(t)=2\eta(t)=2z-\frac{2z^3}{6}+O(z^5)
=\sqrt{t}-\frac{t^{3/2}}{24}+O(t^{5/2}).
\]
并且 \(q(t)=e^{-2\eta(t)}=e^{-f(t)}=1-\sqrt{t}+O(t)\)，故
\[
g(t)=-\log(1+q(t))
=-\log\!\bigl(2-\sqrt{t}+O(t)\bigr)
=-\log2+\frac{\sqrt{t}}{2}+O(t).
\]
分支点位置来自 \(\sqrt{t(4+t)}\) 在 \(t=0,-4\) 的平方根退化。证毕。
\end{proof}

\begin{proposition}[边界反射系数的 Schur--Riccati 递推与黄金耦合下的 Fibonacci 收敛率]\label{prop:pom-Lk-boundary-riccati-recursion}
定义有限尺度边界反射系数
\[
q_k(t):=(L_k+tI)^{-1}_{11}\qquad(t>0).
\]
则 \((q_k(t))_{k\ge 1}\) 满足 Schur--Riccati 递推
\[
q_{k+1}(t)=\frac{1}{t+2-q_k(t)},\qquad q_1(t)=\frac{1}{1+t},
\]
并收敛到固定点 \(q(t)\in(0,1)\)（推论 \ref{cor:pom-Lk-boundary-fixed-point-riccati}）：
\[
q(t)=\frac{1}{t+2-q(t)}\qquad\Longleftrightarrow\qquad q(t)^2-(2+t)q(t)+1=0.
\]
更强地，存在精确误差形式
\[
0<q_k(t)-q(t)=\frac{(1-q(t)^2)\,q(t)^{2k}}{1+q(t)^{2k+1}}
< (1-q(t)^2)\,q(t)^{2k},
\]
从而 \(|q_k(t)-q(t)|=\Theta(q(t)^{2k})\)。
在黄金耦合 \(t=1\) 下，该递推完全算术化：
\[
q_k(1)=\frac{F_{2k-1}}{F_{2k+1}}\ \xrightarrow[k\to\infty]{}\ \varphi^{-2},
\qquad
0<q_k(1)-\varphi^{-2}=\Theta(\varphi^{-4k}).
\]
\leanverified{pom_Lk_boundary_riccati_recursion_t1}
\leanverified{paper_pom_Lk_boundary_riccati_recursion_package}
\end{proposition}
\begin{proof}
递推由对称三对角矩阵 \(L_k+tI\) 的逐步 Schur 补（等价于一维 continued fraction）得到；也可由命题 \ref{prop:pom-Lk-green-kernel-entries} 的闭式 \(q_k(t)=\cosh((2k-1)\eta(t))/\cosh((2k+1)\eta(t))\) 直接验证。

对误差，令 \(q=q(t)=e^{-2\eta(t)}\)（推论 \ref{cor:pom-Lk-boundary-fixed-point-riccati}），则
\[
q_k(t)=\frac{\cosh((2k-1)\eta)}{\cosh((2k+1)\eta)}
=q\,\frac{1+q^{2k-1}}{1+q^{2k+1}},
\]
从而
\[
q_k(t)-q=q\,\frac{q^{2k-1}-q^{2k+1}}{1+q^{2k+1}}
=\frac{(1-q^2)\,q^{2k}}{1+q^{2k+1}},
\]
并立得所示双边界与 \(\Theta(q^{2k})\)。
最后，黄金耦合的 Fibonacci 表达由推论 \ref{cor:pom-Lk-t1-fibonacci-det-green} 的 \((A_k^{-1})_{11}=F_{2k-1}/F_{2k+1}\) 给出，且 \(\varphi^{-2}=q(1)\)。
由 Binet 公式可得误差阶 \(\Theta(\varphi^{-4k})\)。证毕。
\end{proof}

\begin{corollary}[黄金耦合 \texorpdfstring{$t=1$}{t=1} 的边界反射误差闭式]\label{cor:pom-Lk-t1-error-closed-form}
在黄金耦合 \(t=1\) 下，令 \(q(1)=\varphi^{-2}\)。
则对一切 \(k\ge 1\) 有完全显式闭式
\begin{equation}\label{eq:pom-Lk-t1-error-closed-form}
\boxed{\;
q_k(1)-\varphi^{-2}
=(1-\varphi^{-4})\frac{\varphi^{-4k}}{1+\varphi^{-(4k+2)}}
=\frac{\varphi^4-1}{\varphi^{4k+4}+\varphi^{2}}\; }.
\end{equation}
从而
\[
0<q_k(1)-\varphi^{-2}<(1-\varphi^{-4})\varphi^{-4k},
\qquad
|q_k(1)-\varphi^{-2}|\asymp \varphi^{-4k}.
\]
\end{corollary}
\leanverified{paper\_pom\_Lk\_t1\_error\_closed\_form}
\leanverified{paper\_pom\_lk\_t1\_error\_closed\_form}
\begin{proof}
由命题 \ref{prop:pom-Lk-boundary-riccati-recursion} 的精确误差式
\(
q_k(t)-q(t)=\frac{(1-q(t)^2)q(t)^{2k}}{1+q(t)^{2k+1}}
\)
并取 \(t=1\) 与 \(q(1)=\varphi^{-2}\)，得
\[
q_k(1)-\varphi^{-2}
=(1-\varphi^{-4})\frac{\varphi^{-4k}}{1+\varphi^{-(4k+2)}}.
\]
将分子分母同乘 \(\varphi^{4k+4}\) 即得第二个闭式表达。
上界由 \(1+\varphi^{-(4k+2)}>1\) 直接推出；
渐近 \(\asymp\varphi^{-4k}\) 由 \(\varphi^{-(4k+2)}\to 0\) 得到。证毕。
\end{proof}

\begin{corollary}[黄金层级截断误差的绝对可求和性]\label{cor:pom-Lk-t1-error-summable}
令 \(\Delta_k:=q_k(1)-\varphi^{-2}\)。
则
\[
0<\Delta_k<(1-\varphi^{-4})\varphi^{-4k},
\qquad
\sum_{k\ge 1}\Delta_k<\infty.
\]
\leanverified{riccati\_error\_pos}
\leanverified{riccati\_error\_lt\_geom}
\leanverified{paper\_pom\_Lk\_t1\_error\_summable}
\end{corollary}
\begin{proof}
由推论 \ref{cor:pom-Lk-t1-error-closed-form} 的上界与几何级数
\(
\sum_{k\ge 1}\varphi^{-4k}<\infty
\)
立即得到可求和性。证毕。
\end{proof}

\begin{corollary}[Riccati 近似的精确误差分解：由 \texorpdfstring{$D_k$}{Dk} 控制的单调收敛与级数余项]\label{cor:pom-Lk-riccati-error-by-determinants}
对 \(k\ge 0\) 定义 \(D_k(t):=\det(L_k+tI)\) 并约定 \(D_0(t)=1\)。
则对一切 \(k\ge 1\) 有精确恒等式
\[
\boxed{\;
q_k(t)=\frac{D_{k-1}(t)}{D_k(t)}\; }.
\]
并且由命题 \ref{prop:pom-Lk-det-cassini-pell} 得差分公式
\[
\boxed{\;
q_k(t)-q_{k+1}(t)=\frac{t}{D_k(t)\,D_{k+1}(t)}\; }.
\]
因此当 \(t>0\) 时，\((q_k(t))_{k\ge 1}\) 严格单调递减并收敛到固定点 \(q(t)\in(0,1)\)，且误差具有完全显式的级数余项表示
\[
\boxed{\;
q_k(t)-q(t)=t\sum_{n\ge k}\frac{1}{D_n(t)\,D_{n+1}(t)}\; }.
\]
\leanverified{paper\_pom\_Lk\_riccati\_error\_by\_determinants}
\end{corollary}

\begin{proof}
令 \(A_k:=L_k+tI\)。
由 Cramer 公式，
\(
(A_k^{-1})_{11}=\det(A_k^{(1,1)})/\det(A_k)
\)，其中 \(A_k^{(1,1)}\) 为删去第 \(1\) 行第 \(1\) 列所得主子式。
由于 \(A_k^{(1,1)}\) 与 \(L_{k-1}+tI\) 在显式三对角结构上同构，故 \(\det(A_k^{(1,1)})=D_{k-1}(t)\)，从而 \(q_k=D_{k-1}/D_k\)。
差分公式由
\[
q_k-q_{k+1}
=\frac{D_{k-1}}{D_k}-\frac{D_k}{D_{k+1}}
=\frac{D_{k-1}D_{k+1}-D_k^2}{D_kD_{k+1}}
=\frac{t}{D_kD_{k+1}}
\]
（命题 \ref{prop:pom-Lk-det-cassini-pell}）立得。
对 \(t>0\) 有 \(D_k(t)>0\)，故 \(q_k-q_{k+1}>0\) 并给出单调性；余项表示由对差分公式从 \(n=k\) 到 \(\infty\) 求和并用 \(q_n\to q\) 得到。证毕。
\end{proof}

\begin{remark}[有限窗审计：黄金耦合 Fibonacci 化、Fisher 零点与 Riccati 收敛]\label{rem:pom-fence-green-kernel-golden-coupling-audit}
脚本在有限窗上核验：推论 \ref{cor:pom-Lk-golden-coupling-unique} 的 \(t_\star=1\) 唯一性刻画、
推论 \ref{cor:pom-Lk-t1-fibonacci-det-green} 的 Fibonacci 行列式与矩阵元闭式、
定理 \ref{thm:pom-Lk-fisher-zeros-arcsine} 的零点位置与反正弦律（通过映射 \(t=-\mu\) 与 Joukowsky 均匀化的数值一致性），
以及命题 \ref{prop:pom-Lk-boundary-riccati-recursion} 的 Riccati 递推与收敛率。
审计表如下：
\begin{table}[H]
\centering
\scriptsize
\setlength{\tabcolsep}{6pt}
\caption{黄金耦合与 Fisher 零点的有限窗审计：$\det(L_k+I)=F_{2k+1}$ 与 Fibonacci 矩阵元闭式（推论~\ref{cor:pom-Lk-t1-fibonacci-det-green}）、零点反正弦律（定理~\ref{thm:pom-Lk-fisher-zeros-arcsine}）、Joukowsky 均匀化的奇根条件（注~\ref{rem:pom-Lk-fisher-zeros-uniformization}）、以及 Riccati 递推的对象层一致性（命题~\ref{prop:pom-Lk-boundary-riccati-recursion}）。}
\label{tab:fence_green_kernel_golden_coupling_audit}
\begin{tabular}{r c r r r r}
\toprule
$k$ & det ok & $\max|\Delta A^{-1}_{ij}|$ & $\|F_k-F\|_\infty$ & $|\Delta q_k|$ & $\max|w^{2k+1}+1|$\\
\midrule
10 & yes & 2.78e-17 & 9.52e-02 & 0.00e+00 & 8.30e-15\\
20 & yes & 2.78e-17 & 4.88e-02 & 0.00e+00 & 1.21e-14\\
40 & yes & 2.78e-17 & 2.47e-02 & 0.00e+00 & 2.45e-14\\
200 & yes & 0.00e+00 & 4.99e-03 & 0.00e+00 & 1.90e-13\\
\bottomrule
\end{tabular}
\end{table}

\end{remark}

\paragraph{\texorpdfstring{$k$}{k}-重谱碰撞的根单位过滤核：递推、临界缩放与零点指纹}
为刻画高重并根临界点附近的统一核结构，固定整数 \(k\ge 2\)，定义
\[
Z_n^{(k)}(t):=\sum_{r\ge 0}\binom{n+k-1}{k-1+kr}\,t^r,\qquad n\in\ZZ_{\ge 0}.
\]
其中求和在 \(k-1+kr\le n+k-1\) 时终止。

\begin{theorem}[根单位过滤配分函数的代数结构]\label{thm:pom-kcollision-root-filter-recurrence}
对上述 \(Z_n^{(k)}(t)\) 有：
\begin{enumerate}
\item 系数均为非负整数，且 \(\deg_t Z_n^{(k)}=\lfloor n/k\rfloor\)；
\item 序列 \((Z_n^{(k)}(t))_{n\ge 0}\) 满足阶 \(k\) 常系数递推
\[
Z_{n+k}^{(k)}
=\binom{k}{1}Z_{n+k-1}^{(k)}-\binom{k}{2}Z_{n+k-2}^{(k)}+\cdots+(-1)^{k-1}\binom{k}{k-1}Z_{n+1}^{(k)}
+\bigl(t-(-1)^k\bigr)Z_n^{(k)},
\]
且在 \(0\le n\le k-1\) 时
\(
Z_n^{(k)}(t)=\binom{n+k-1}{k-1}
\)；
\item 对应特征多项式为
\[
\chi_k(\lambda;t)=(\lambda-1)^k-t.
\]
\end{enumerate}
\leanverified{paper\_pom\_kcollision\_root\_filter\_recurrence\_seeds}
\leanverified{paper\_pom\_kcollision\_root\_filter\_recurrence\_package}
\leanverified{paper\_pom\_kcollision\_root\_filter\_extended}
\leanverified{falling\_factorial\_seeds}
\leanverified{descFactorial\_eq\_factorial\_mul\_choose}
\leanverified{stirling\_second\_kind\_seeds}
\end{theorem}
\begin{proof}
第一点由定义直接成立。
将 \((\lambda-1)^k-t\) 展开为
\[
\lambda^k-\binom{k}{1}\lambda^{k-1}+\binom{k}{2}\lambda^{k-2}-\cdots+(-1)^k-t
\]
并代入常系数递推与特征多项式的对应关系，得到第二、三点。
初值由 \(\lfloor n/k\rfloor=0\) 时仅 \(r=0\) 项贡献给出。证毕。
\end{proof}

\begin{theorem}[重碰撞临界窗口的 Mittag--Leffler 极限]\label{thm:pom-kcollision-mittag-leffler-scaling}
令
\[
E_{\alpha,\beta}(z):=\sum_{m=0}^{\infty}\frac{z^m}{\Gamma(\alpha m+\beta)},\qquad \alpha>0.
\]
固定 \(k\ge 2\)。当 \(n\to\infty\) 时，在任意紧集 \(K\subset\CC\) 上一致有
\[
n^{-(k-1)}Z_n^{(k)}\!\left(-\frac{u}{n^k}\right)\longrightarrow E_{k,k}(-u),\qquad u\in K.
\]
\end{theorem}
\begin{proof}
写作
\[
Z_n^{(k)}\!\left(-\frac{u}{n^k}\right)
=\sum_{r\ge 0}\binom{n+k-1}{k-1+kr}\left(-\frac{u}{n^k}\right)^r.
\]
固定 \(r\) 时，
\[
\binom{n+k-1}{k-1+kr}
=\frac{n^{k-1+kr}}{(k-1+kr)!}\bigl(1+o(1)\bigr),
\]
故
\[
n^{-(k-1)}\binom{n+k-1}{k-1+kr}\left(-\frac{u}{n^k}\right)^r
\to \frac{(-u)^r}{\Gamma(kr+k)}.
\]
再由
\[
\left|n^{-(k-1)}\binom{n+k-1}{k-1+kr}\left(\frac{u}{n^k}\right)^r\right|
\le \frac{C_K^r\,e^{kr}}{(k-1+kr)!},\qquad u\in K,
\]
及右端可求和，应用 Weierstrass \(M\)-判别法可交换极限与求和，得结论。证毕。
\end{proof}
\leanverified{paper\_pom\_kcollision\_mittag\_leffler\_scaling}

\begin{corollary}[Lee--Yang 零点的临界缩放极限]\label{cor:pom-kcollision-lee-yang-scaling}
记 \(E_{k,k}(-u)\) 的零点（按重数）为 \(\{\xi_{k,j}\}_{j\ge 1}\subset\CC\setminus\{0\}\)。
对任意固定 \(J\in\NN\)，当 \(n\) 充分大时，\(Z_n^{(k)}(t)\) 在窗口 \(|t|=O(n^{-k})\) 内对应的前 \(J\) 个零点满足
\[
t_{n,j}\sim -\frac{\xi_{k,j}}{n^k},\qquad 1\le j\le J.
\]
并且
\[
\{-n^k t:\ Z_n^{(k)}(t)=0\}\xrightarrow[n\to\infty]{\mathrm{loc.\ Hausdorff}}\{\xi_{k,j}\}_{j\ge 1}.
\]
\end{corollary}
\leanverified{paper\_pom\_kcollision\_lee\_yang\_scaling}
\begin{proof}
令
\[
F_n(u):=n^{-(k-1)}Z_n^{(k)}\!\left(-\frac{u}{n^k}\right),\qquad F(u):=E_{k,k}(-u).
\]
由定理 \ref{thm:pom-kcollision-mittag-leffler-scaling}，\(F_n\to F\) 在紧集上一致，且均全纯。
由 Hurwitz 定理，紧域内零点（计重数）收敛到 \(F\) 的零点。
变量替换 \(t=-u/n^k\) 即得。证毕。
\end{proof}

\begin{proposition}[Hankel 行列式的秩截断与主指纹]\label{prop:pom-kcollision-hankel-fingerprint}
定义
\[
H_m^{(k)}(t):=\det\bigl[Z_{i+j}^{(k)}(t)\bigr]_{0\le i,j\le m-1}.
\]
则
\begin{enumerate}
\item 若 \(m>k\)，则 \(H_m^{(k)}(t)\equiv 0\)；
\item 主 Hankel 行列式满足
\[
H_k^{(k)}(t)=(-1)^{\frac{(k-1)(3k-4)}{2}}\bigl(1-(-1)^k t\bigr)^{k-1}.
\]
\end{enumerate}
\leanverified{paper\_pom\_kcollision\_hankel\_fingerprint\_seeds}
\leanverified{paper\_pom\_kcollision\_hankel\_fingerprint\_package}
\end{proposition}
\begin{proof}
取 \(\tau^k=t\)、\(\omega_j=e^{2\pi i j/k}\)。
由根单位过滤恒等式，
\[
Z_n^{(k)}(t)=\sum_{j=0}^{k-1}w_j(t)\lambda_j(t)^n,\qquad
\lambda_j(t):=1+\omega_j\tau,\quad
w_j(t):=\frac{1}{k}\frac{\lambda_j(t)^{k-1}}{\omega_j^{k-1}\tau^{k-1}}.
\]
故 Hankel 矩阵分解为 \(V_m^\top W V_m\)，其中
\[
(V_m)_{j,i}=\lambda_j(t)^i,\qquad
W=\mathrm{diag}(w_0,\dots,w_{k-1}).
\]
从而 \(\mathrm{rank}\le k\)，得 \(m>k\) 时行列式为零。
当 \(m=k\) 时，
\[
H_k^{(k)}=(\det V_k)^2\prod_{j=0}^{k-1}w_j.
\]
再用 Vandermonde 公式、\(x^k-1\) 判别式
\[
\prod_{0\le i<j\le k-1}(\omega_j-\omega_i)^2
=(-1)^{\frac{(k-1)(k-2)}{2}}k^k,
\]
以及
\[
\prod_{j=0}^{k-1}(1+\omega_j\tau)=1-(-1)^k t,\qquad
\prod_{j=0}^{k-1}\omega_j=(-1)^{k-1},
\]
整理即得所示闭式。证毕。
\end{proof}

\begin{theorem}[\texorpdfstring{$E_{k,k}$}{Ekk} 的根单位指数分解与实轴零点量子化]\label{thm:pom-kcollision-ekk-root-phase-quantization}
对任意分支 \(z^{1/k}\) 有恒等式
\[
E_{k,k}(z)=\frac1k\,z^{\frac{1-k}{k}}\sum_{j=0}^{k-1}\omega_j^{1-k}\exp\!\bigl(\omega_j z^{1/k}\bigr),\qquad
\omega_j=e^{2\pi i j/k}.
\]
令 \(u>0\)，取主值分支 \((-u)^{1/k}=u^{1/k}e^{i\pi/k}\)。
当 \(u\to+\infty\) 时，
\[
E_{k,k}(-u)
=-\frac{2}{k}u^{\frac{1-k}{k}}
\exp\!\Bigl(u^{1/k}\cos\frac{\pi}{k}\Bigr)
\cos\!\Bigl(u^{1/k}\sin\frac{\pi}{k}+\frac{\pi}{k}\Bigr)
+O\!\left(u^{\frac{1-k}{k}}\exp\!\Bigl(u^{1/k}\cos\frac{3\pi}{k}\Bigr)\right),
\]
其中 \(k=2\) 时误差项消失。
进一步，记 \(E_{k,k}(-u)\) 的正实大零点为 \(\xi_{k,n}\)（按递增计），则
\[
\xi_{k,n}
=\left(\frac{\bigl(n+\frac12-\frac1k\bigr)\pi}{\sin(\pi/k)}\right)^{k}\bigl(1+o(1)\bigr),\qquad n\to\infty.
\]
\leanverified{paper\_pom\_kcollision\_ekk\_root\_phase\_quantization}
\end{theorem}
\begin{proof}
第一式由
\(
e^{\omega_j z^{1/k}}=\sum_{m\ge 0}\omega_j^m z^{m/k}/m!
\)
与根单位正交关系筛选 \(m\equiv k-1\pmod k\) 得到。
第二式中，代入 \(z=-u\) 后，\(j=0\) 与 \(j=k-1\) 两项给出最大实部并构成共轭主导对，其余项实部不超过 \(\cos(3\pi/k)\) 对应指数级次，故得主项加误差。
令主项余弦相位等于 \((n+\frac12)\pi\) 即得零点量子化近似；指数分离保证相对误差为 \(o(1)\)。证毕。
\end{proof}

\begin{conjecture}[有限维转移模型中的重碰撞临界核普适性]\label{conj:pom-kcollision-universality-ekk}
设 \((Z_n(t))\) 来自有限维转移矩阵（等价常系数递推）并在 \(t=t_\ast\) 发生原始 \(k\)-重主导谱碰撞：主导特征值代数重数为 \(k\)，其余谱与主导谱存在正模长间隙，且无额外对称性引起进一步约化。
则存在常数 \(C,c\ne 0\) 使在临界窗口 \(t=t_\ast-u/n^k\) 下
\[
n^{-(k-1)}Z_n\!\left(t_\ast-\frac{u}{n^k}\right)\Longrightarrow C\,E_{k,k}(-c\,u),
\]
并且 Lee--Yang 零点在该窗口内按推论 \ref{cor:pom-kcollision-lee-yang-scaling} 收敛到相应 \(E_{k,k}\) 零点集。
\end{conjecture}


\begin{corollary}[谱隙常数的完全显式化]\label{cor:pom-Lk-gap-pi2}
\leanverified{paper\_pom\_lk\_gap\_pi2}
令 \(\mu_1\le\mu_2\le\cdots\le\mu_k\) 为 \(L_k\) 的升序特征值，则
\[
\mu_1=4\sin^2\!\Bigl(\frac{\pi}{4k+2}\Bigr),
\qquad
\mu_2=4\sin^2\!\Bigl(\frac{3\pi}{4k+2}\Bigr),
\qquad
\frac{\mu_2}{\mu_1}=\frac{\sin^2\!\bigl(\frac{3\pi}{4k+2}\bigr)}{\sin^2\!\bigl(\frac{\pi}{4k+2}\bigr)}.
\]
并且当 \(k\to\infty\) 时有渐近
\[
\mu_1=\frac{\pi^2}{(2k+1)^2}+O(k^{-4}),
\qquad
\lambda_{\max}(K_k)=\mu_1^{-1}=\frac{(2k+1)^2}{\pi^2}+O(1).
\]
\end{corollary}
\begin{proof}
由注 \ref{rem:pom-Kk-factor4-spectral-rigidity} 的谱闭式与 \(\phi_p=\frac{(2p-1)\pi}{4k+2}\) 的单调性可得最小/次小特征值表达。
最后两式由 \(\sin x=x+O(x^3)\) 得 \(\mu_1=4(\pi/(4k+2))^2+O(k^{-4})\)，并取倒数得到 \(\lambda_{\max}(K_k)=\mu_1^{-1}\)。证毕。
\end{proof}

\begin{proposition}[Green 核 \texorpdfstring{$K_k$}{Kk} 的偶次多项式矩律诱导 \texorpdfstring{$\zeta(2r)$}{zeta(2r)} 极限：Bernoulli 系数与奇整数谱矩锁合（228）]\label{prop:pom-csc-moments-bernoulli}
对整数 \(r\ge 1\) 定义
\[
S_r(k):=\sum_{p=1}^{k}\csc^{2r}\!\Bigl(\frac{(2p-1)\pi}{4k+2}\Bigr),
\qquad
n:=2k+1.
\]
则 \(S_r(k)\) 是 \(n\) 的偶次多项式，且存在有理系数 \(a_{r,j}\in\QQ\) 使
\[
S_r(k)=-\frac12+\sum_{j=1}^{r} a_{r,j}\,n^{2j}.
\]
其最高次系数由 Bernoulli 数给出：
\[
a_{r,r}=\frac{2^{2r-1}(2^{2r}-1)\,|B_{2r}|}{(2r)!}.
\]
并且 \(K_k\) 的迹矩满足
\[
\mathrm{tr}(K_k^{\,r})=\sum_{p=1}^{k}\lambda_p^{\,r}=\frac{1}{4^{\,r}}\,S_r(k),
\]
其中 \(\lambda_p=(4\sin^2\phi_p)^{-1}\) 为 \(K_k\) 的特征值（引理 \ref{lem:pom-Kk-eigenvalues}）。
特别地，取 \(r\in\NN\) 并按主尺度 \(n=2k+1\) 归一化，则由推论 \ref{cor:pom-Lk-normalized-zeta-odd-projection}（注意 \(\zeta_k(r)=\sum_p\mu_p^{-r}=\sum_p\lambda_p^{\,r}=\mathrm{tr}(K_k^{\,r})\)）有严格极限
\[
\lim_{k\to\infty}\Bigl(\frac{\pi}{n}\Bigr)^{2r}\mathrm{tr}(K_k^{\,r})
=(1-2^{-2r})\zeta(2r)
=\sum_{\substack{q\ge 1\\ q\ \mathrm{odd}}}\frac{1}{q^{2r}}.
\]
\end{proposition}
\begin{proof}
令
\(
A_r(m):=\sum_{q=1}^{m-1}\csc^{2r}(\pi q/m)
\)。
由 \(\pi\cot(\pi z)\) 的部分分式展开与留数计算可得 \(A_r(m)\) 为 \(m\) 的偶次多项式，且其最高次系数为 \(2^{2r}|B_{2r}|/(2r)!\)。
将 \(m=2n\) 的总和按奇偶指标分解：
\[
A_r(2n)=\sum_{q=1}^{n-1}\csc^{2r}\!\Bigl(\frac{\pi q}{n}\Bigr)+\sum_{\substack{1\le q\le 2n-1\\ q\ \mathrm{odd}}}\csc^{2r}\!\Bigl(\frac{\pi q}{2n}\Bigr)
=A_r(n)+\Bigl(2S_r(k)+1\Bigr),
\]
其中最后一步用到正弦对称性把奇指标和配对为 \(2S_r(k)\)，并分离出 \(q=n\) 的 \(\csc^{2r}(\pi/2)=1\)。
故
\(
S_r(k)=\frac12(A_r(2n)-A_r(n)-1)
\)，从而 \(S_r(k)\) 为 \(n\) 的偶次多项式并具有常数项 \(-1/2\)。
记 \(c\) 为 \(A_r(m)\) 的最高次系数，则 \(A_r(2n)\) 的 \(n^{2r}\) 系数为 \(2^{2r}c\)，而 \(A_r(n)\) 的 \(n^{2r}\) 系数为 \(c\)，故 \(S_r(k)\) 的最高次系数为
\[
\frac{2^{2r}-1}{2}\,c.
\]
代入 \(c=2^{2r}|B_{2r}|/(2r)!\) 即得所示
\(
a_{r,r}=\frac{2^{2r-1}(2^{2r}-1)\,|B_{2r}|}{(2r)!}
\)。
最后一式由 \(\lambda_p=(4\sin^2\phi_p)^{-1}=(1/4)\csc^2\phi_p\) 直接代入即得。证毕。
\end{proof}

\leanverified{paper\_pom\_csc\_moments\_finite\_trace\_inversion}

\begin{corollary}[偶值 \texorpdfstring{$\zeta$}{zeta} 的有限迹反演：由有限个离散迹矩确定主项系数]\label{cor:pom-csc-moments-finite-trace-inversion}
固定整数 \(r\ge 1\)，并对 \(k\ge 1\) 记 \(n_k:=2k+1\)、\(u_k:=n_k^2\)。
令 \(S_r(k)\) 如命题 \ref{prop:pom-csc-moments-bernoulli}，并写作
\[
S_r(k)+\frac12=\sum_{j=1}^{r}a_{r,j}\,n_k^{2j}=:Q_r(u_k),
\qquad
Q_r(U):=\sum_{j=1}^{r}a_{r,j}\,U^{j}\in\QQ[U].
\]
则 \(Q_r\) 为次数 \(r\) 的多项式，并由 \(r+1\) 个有限迹矩
\(
\{\mathrm{tr}(K_k^{\,r})\}_{k=1}^{r+1}
\)
（等价地由 \(\{S_r(k)\}_{k=1}^{r+1}\)）唯一确定；
从而主项系数 \(a_{r,r}\) 可由有限数据确定性回收，并满足恒等式
\[
\boxed{\;
(1-2^{-2r})\zeta(2r)=\frac{\pi^{2r}}{4^{\,r}}\,a_{r,r}\; }.
\]
\end{corollary}
\begin{proof}
由命题 \ref{prop:pom-csc-moments-bernoulli}，对每个 \(k\) 有 \(S_r(k)=-\frac12+Q_r(n_k^2)\) 且 \(Q_r\in\QQ[U]\) 的次数为 \(r\)。
又 \(u_k=n_k^2\) 在 \(k=1,\dots,r+1\) 上互异，故插值唯一性推出 \(Q_r\)（因而各 \(a_{r,j}\)）由 \(\{Q_r(u_k)\}_{k=1}^{r+1}\) 唯一确定。
而 \(Q_r(u_k)=S_r(k)+\frac12=4^{\,r}\mathrm{tr}(K_k^{\,r})+\frac12\) 由有限迹矩直接给出。
最后，\(\mathrm{tr}(K_k^{\,r})=\frac{1}{4^{\,r}}S_r(k)=\frac{a_{r,r}}{4^{\,r}}\,n_k^{2r}+O(n_k^{2r-2})\)，故
\(
(\pi/n_k)^{2r}\mathrm{tr}(K_k^{\,r})\to \pi^{2r}a_{r,r}/4^{\,r}
\)；
再与命题 \ref{prop:pom-csc-moments-bernoulli} 的极限恒等式对齐即得。证毕。
\end{proof}

\begin{remark}[Green 核 \texorpdfstring{$K_k$}{Kk} 的前三阶矩闭式与 Schatten 标度：有限窗即可完全审计（229）]\label{rem:pom-csc-moments-examples}
写 \(n:=2k+1\)。
则
\[
S_1(k)=\frac{n^2}{2}-\frac12,\qquad
S_2(k)=\frac{n^4}{6}+\frac{n^2}{3}-\frac12,\qquad
S_3(k)=\frac{n^6}{15}+\frac{n^4}{6}+\frac{4n^2}{15}-\frac12.
\]
从而 \(\mathrm{tr}(K_k^{\,r})=\frac{1}{4^{\,r}}S_r(k)\) 在 \(r=1,2,3\) 的有限窗上可逐项核验（表 \ref{tab:fence_green_kernel_spectral_algebra_audit}）。
由于 \(K_k\succeq 0\)，上述闭式亦为 \(\|K_k\|_{S_r}^r=\mathrm{tr}(K_k^{\,r})\)（\(r=1,2,3\)）提供无需数值对角化的确定型输入。
\end{remark}

\begin{remark}[有限窗审计：Gram/行列式刚性、Chebyshev 识别与高阶矩]\label{rem:pom-fence-green-kernel-spectral-algebra-audit}
脚本在有限窗上逐项核验：\(\det(K_k)=1\) 及由此推出的 \((4\sin^2)\) 乘积/求和恒等式（推论 \ref{cor:pom-Kk-sine-product-sum}）、
Chebyshev 特征多项式识别（定理 \ref{thm:pom-Lk-chebyshev-charpoly}）、
以及前三阶幂和闭式（注 \ref{rem:pom-csc-moments-examples}）。
审计表如下：
\begin{table}[H]
\centering
\scriptsize
\setlength{\tabcolsep}{6pt}
\caption{Green 核 $K_k(i,j)=\min(i,j)$ 的谱代数审计：核验 Gram 分解与 $\det(K_k)=1$ 的行列式归一化刚性（引理~\ref{lem:pom-Kk-gram-det}），$(4\sin^2)$ 乘积/求和恒等式（推论~\ref{cor:pom-Kk-sine-product-sum}），混合边界 Laplacian $L_k=K_k^{-1}$ 的 Chebyshev 特征多项式识别（定理~\ref{thm:pom-Lk-chebyshev-charpoly}），以及移位行列式自由能重整化常数项与谱隙尺度（推论~\ref{cor:pom-Lk-shifted-det-free-energy}、\ref{cor:pom-Lk-gap-pi2}）。}
\label{tab:fence_green_kernel_spectral_algebra_audit}
\begin{tabular}{r r r c c c r r}
\toprule
$k$ & $|\prod 4\sin^2-1|$ & $|\sum (4\sin^2)^{-1}-\tfrac{k(k+1)}{2}|$ & Cheb & $\det$-shift & moments & $\mu_1(2k+1)^2/\pi^2$ & $|F_k(1)-F_\infty(1)|$\\
\midrule
1 & 2.22e-16 & 2.22e-16 & OK & OK & OK & 0.91189065 & 5.42e-02\\
2 & 1.11e-16 & 4.44e-16 & OK & OK & OK & 0.96753121 & 8.10e-03\\
3 & 3.33e-16 & 0.00e+00 & OK & OK & OK & 0.98332725 & 1.19e-03\\
4 & 5.55e-16 & 1.78e-15 & OK & OK & OK & 0.98988724 & 1.73e-04\\
5 & 5.55e-16 & 0.00e+00 & OK & OK & OK & 0.99322121 & 2.53e-05\\
6 & 2.22e-16 & 3.55e-15 & OK & OK & OK & 0.99514280 & 3.68e-06\\
7 & 5.55e-16 & 1.07e-14 & OK & OK & OK & 0.99634993 & 5.37e-07\\
8 & 1.55e-15 & 0.00e+00 & -- & -- & OK & 0.99715733 & 7.84e-08\\
9 & 4.44e-16 & 0.00e+00 & -- & -- & OK & 0.99772377 & 1.14e-08\\
10 & 7.77e-16 & 7.11e-15 & -- & -- & OK & 0.99813639 & 1.67e-09\\
11 & 3.33e-16 & 0.00e+00 & -- & -- & OK & 0.99844621 & 2.44e-10\\
12 & 9.99e-16 & 1.42e-14 & -- & -- & OK & 0.99868475 & 3.55e-11\\
\bottomrule
\end{tabular}
\end{table}

\end{remark}

\paragraph{谱测度与连续极限：反正弦律、离散 Weyl 定律与热核/Bessel 对偶}
下面把混合边界离散 Laplacian \(L_k\) 的“半奇整数量子化”进一步提升为一个\emph{谱测度极限律}，
并给出全区间上的离散 Weyl 计数余项界与热核密度的闭式极限。

\begin{theorem}[经验谱测度的反正弦极限律]\label{thm:pom-Lk-arcsine-law}
\leanverified{paper\_pom\_Lk\_arcsine\_law}
设 \(L_k:=K_k^{-1}\)（见注 \ref{rem:pom-Kk-factor4-spectral-rigidity}），其特征值为
\[
\mu_p
=4\sin^2\!\Bigl(\frac{(2p-1)\pi}{4k+2}\Bigr)
=2\Bigl(1-\cos\frac{(2p-1)\pi}{2k+1}\Bigr)\in(0,4),
\qquad p=1,\dots,k.
\]
定义经验谱测度
\[
\nu_k:=\frac{1}{k}\sum_{p=1}^{k}\delta_{\mu_p}.
\]
则 \(\nu_k\) 在弱拓扑下收敛到 \([0,4]\) 上的反正弦分布 \(\nu\)，其密度为
\[
\dd\nu(\mu)=\frac{1}{\pi\sqrt{\mu(4-\mu)}}\,\dd\mu,\qquad \mu\in(0,4).
\]
等价地，对任意连续函数 \(f\in C([0,4])\)，有迹函数极限
\[
\lim_{k\to\infty}\frac{1}{k}\operatorname{tr}f(L_k)
=\frac{1}{\pi}\int_{0}^{4}\frac{f(\mu)}{\sqrt{\mu(4-\mu)}}\,\dd\mu
=\frac{1}{\pi}\int_{0}^{\pi} f\!\bigl(2(1-\cos\theta)\bigr)\,\dd\theta.
\]
\end{theorem}
\begin{proof}
令 \(n:=2k+1\) 并记 \(\theta_p:=\frac{(2p-1)\pi}{n}\in(0,\pi)\)，则 \(\mu_p=2(1-\cos\theta_p)\)。
因为 \(\{\theta_p\}_{p=1}^k\) 为 \((0,\pi)\) 上步长 \(2\pi/n\) 的半奇整数网格，
对任意连续 \(g\in C([0,\pi])\) 有 Riemann 和收敛
\[
\lim_{k\to\infty}\frac{1}{k}\sum_{p=1}^{k} g(\theta_p)
=\frac{1}{\pi}\int_0^\pi g(\theta)\,\dd\theta.
\]
取 \(g(\theta):=f(2(1-\cos\theta))\) 即得
\(
\lim_{k\to\infty}\frac{1}{k}\sum_{p} f(\mu_p)=\frac{1}{\pi}\int_0^\pi f(2(1-\cos\theta))\,\dd\theta
\)。
再作代换 \(\mu=2(1-\cos\theta)\)（于是 \(\mu(4-\mu)=4\sin^2\theta\) 且 \(\dd\mu=2\sin\theta\,\dd\theta=\sqrt{\mu(4-\mu)}\,\dd\theta\)），得到
\[
\frac{1}{\pi}\int_0^\pi f(2(1-\cos\theta))\,\dd\theta
=\frac{1}{\pi}\int_0^4 \frac{f(\mu)}{\sqrt{\mu(4-\mu)}}\,\dd\mu,
\]
从而结论成立。证毕。
\end{proof}
% Boundary spectral measure for L_k: exact weights and Beta(3/2,3/2) limit.

\begin{theorem}[边界谱测度的精确权重与 \texorpdfstring{Beta\((\tfrac32,\tfrac32)\)}{Beta(3/2,3/2)} 极限]\label{thm:pom-Lk-boundary-spectral-measure-beta32}
沿用定理 \ref{thm:pom-Lk-arcsine-law} 的特征角与特征值记号
\[
\theta_p=\frac{(2p-1)\pi}{2k+1},
\qquad
\mu_p=2(1-\cos\theta_p)=4\sin^2\!\Bigl(\frac{\theta_p}{2}\Bigr),
\qquad p=1,\dots,k.
\]
以边界向量 \(e_1\in\RR^k\) 定义谱测度
\[
\rho_k:=\sum_{p=1}^{k} w_{p,k}\,\delta_{\mu_p},
\qquad
w_{p,k}:=\bigl|\langle e_1,v_p\rangle\bigr|^2,
\]
其中 \(v_p\) 为 \(L_k\) 的单位特征向量。
则权重具有严格闭式
\[
\boxed{\;
w_{p,k}=\frac{4}{2k+1}\sin^2\theta_p,
\qquad \sum_{p=1}^{k} w_{p,k}=1\; }.
\]
并且 \(\rho_k\) 在弱拓扑下收敛到 \([0,4]\) 上的极限测度 \(\rho\)，其密度为
\[
\boxed{\;
\dd\rho(\mu)=\frac{\sqrt{\mu(4-\mu)}}{2\pi}\,\dd\mu,\qquad \mu\in(0,4)\; }.
\]
等价地，令 \((a)_n\) 为 Pochhammer 符号，则对任意整数 \(n\ge 0\) 有矩闭式
\[
\int_{0}^{4}\mu^{n}\,\dd\rho(\mu)=4^{n}\,\frac{(\tfrac32)_n}{(3)_n}.
\]
\leanverified{paper\_pom\_lk\_boundary\_spectral\_measure\_beta32}
\end{theorem}

\begin{proof}
设 \(\theta\in(0,\pi)\) 且 \(\mu=2(1-\cos\theta)\)。
考虑向量 \(u(\theta)\in\RR^k\) 由
\(
u(\theta)_i:=\sin(i\theta)
\)
（\(1\le i\le k\)）定义。
对 \(1<i<k\) 有差分恒等式
\[
-(u_{i-1})+2u_i-(u_{i+1})
=2(1-\cos\theta)\,u_i=\mu\,u_i.
\]
末端行满足
\[
-(u_{k-1})+u_k
=\bigl(2u_k-u_{k+1}\bigr)-u_{k-1}
=\mu\,u_k\quad\Longleftrightarrow\quad u_{k+1}=u_k.
\]
该边界条件等价于 \(\sin((k+1)\theta)=\sin(k\theta)\)，从而等价于 \(\cos\!\bigl(\tfrac{2k+1}{2}\theta\bigr)=0\)，即
\(\theta=\theta_p=\frac{(2p-1)\pi}{2k+1}\)。
因此 \(u(\theta_p)\) 给出 \(L_k\) 的特征向量，特征值为 \(\mu_p\)。
\par
由半奇整数量子化的离散正交性（离散正弦基的标准恒等式），对 \(p\neq q\) 有
\(
\sum_{i=1}^{k}\sin(i\theta_p)\sin(i\theta_q)=0
\)，并且
\[
\sum_{i=1}^{k}\sin^2(i\theta_p)=\frac{2k+1}{4}.
\]
因此单位特征向量可取为
\[
v_p(i)=\sqrt{\frac{4}{2k+1}}\,\sin(i\theta_p),
\]
从而权重闭式为
\(
w_{p,k}=|v_p(1)|^2=\frac{4}{2k+1}\sin^2\theta_p
\)。
并且 \(\sum_{p}w_{p,k}=\sum_{p}|\langle e_1,v_p\rangle|^2=\|e_1\|^2=1\)。
\par
对任意 \(f\in C([0,4])\)，有
\[
\int f(\mu)\,\dd\rho_k(\mu)
=\frac{4}{2k+1}\sum_{p=1}^{k}\sin^2(\theta_p)\,f\!\bigl(2(1-\cos\theta_p)\bigr).
\]
注意 \(\{\theta_p\}\) 是 \((0,\pi)\) 上步长 \(2\pi/(2k+1)\) 的半奇整数网格。
将上式写成 Riemann 和并取极限，得到
\[
\lim_{k\to\infty}\int f\,\dd\rho_k
=\frac{2}{\pi}\int_0^\pi \sin^2(\theta)\,f\!\bigl(2(1-\cos\theta)\bigr)\,\dd\theta.
\]
作代换 \(\mu=2(1-\cos\theta)\)（此时 \(\dd\mu=\sqrt{\mu(4-\mu)}\,\dd\theta\) 且 \(\sin^2\theta=\mu(4-\mu)/4\)），得到极限测度密度为
\(
\dd\rho(\mu)=\frac{\sqrt{\mu(4-\mu)}}{2\pi}\dd\mu
\)。
矩公式由令 \(\mu=4x\)（\(x\in(0,1)\)）把 \(\rho\) 识别为 \(x\) 的 Beta\((\tfrac32,\tfrac32)\) 分布推前而得。证毕。
\end{proof}

\begin{corollary}[Catalan 矩封闭式与代数生成函数]\label{cor:pom-Lk-boundary-catalan-moments-gf}
令
\(
m_n:=\int_{0}^{4}\mu^{n}\,\dd\rho(\mu)
\)（\(n\ge 0\)）。
则对一切 \(n\ge 0\) 成立严格恒等式
\[
\boxed{\;
m_n=C_{n+1},\qquad
C_n:=\frac{1}{n+1}\binom{2n}{n}\;}
\]
其中 \(C_n\) 为 Catalan 数。
因此矩生成函数
\(
M(z):=\sum_{n\ge 0}m_n z^n
\)
为代数函数并满足闭式
\[
\boxed{\;
M(z)=\sum_{n\ge 0}C_{n+1}z^n
=\frac{1-2z-\sqrt{1-4z}}{2z^2}\; }.
\]
\leanverified{paper\_pom\_catalan\_moments}
\leanverified{paper\_pom\_lk\_boundary\_catalan\_moments\_gf}
\end{corollary}

\begin{proof}
由定理 \ref{thm:pom-Lk-boundary-spectral-measure-beta32}，
\(
m_n=4^{n}\frac{(\tfrac32)_n}{(3)_n}
\)。
利用 \((3)_n=\Gamma(n+3)/\Gamma(3)=(n+2)!/2\) 与 \((\tfrac32)_n=\Gamma(n+\tfrac32)/\Gamma(\tfrac32)=(2n+1)!!/2^{n}\) 得
\[
m_n=4^{n}\cdot\frac{(2n+1)!!}{2^{n}}\cdot\frac{2}{(n+2)!}
=2^{n+1}\frac{(2n+1)!!}{(n+2)!}.
\]
再用恒等式 \((2n+1)!!=\frac{(2n+2)!}{2^{n+1}(n+1)!}\) 化简得到
\[
m_n=\frac{(2n+2)!}{(n+1)!(n+2)!}
=\frac{1}{n+2}\binom{2n+2}{n+1}
=C_{n+1}.
\]
记 Catalan 生成函数 \(C(z)=\sum_{n\ge 0}C_n z^n=\frac{1-\sqrt{1-4z}}{2z}\)，则
\(
M(z)=\sum_{n\ge 0}C_{n+1}z^n=\frac{C(z)-C_0}{z}=\frac{C(z)-1}{z}
\)，代入即得所示闭式。证毕。
\end{proof}

\leanverified{paper\_pom\_lk\_boundary\_bulk\_radon\_nikodym}
\begin{corollary}[边界—体相测度的 Radon--Nikodym 精确关系]\label{cor:pom-Lk-boundary-bulk-radon-nikodym}
记定理 \ref{thm:pom-Lk-arcsine-law} 的体相极限测度为 \(\nu\)。
则定理 \ref{thm:pom-Lk-boundary-spectral-measure-beta32} 的边界极限测度 \(\rho\) 满足精确绝对连续关系
\[
\boxed{\;
\dd\rho(\mu)=\frac{\mu(4-\mu)}{2}\,\dd\nu(\mu)\; }.
\]
因此对任意可积函数 \(f\) 有严格换测公式
\[
\int f(\mu)\,\dd\rho(\mu)=\frac12\int \mu(4-\mu)\,f(\mu)\,\dd\nu(\mu).
\]
\end{corollary}

\begin{proof}
由两测度密度的比值
\(
\frac{\sqrt{\mu(4-\mu)}/(2\pi)}{1/(\pi\sqrt{\mu(4-\mu)})}=\mu(4-\mu)/2
\)
立得。证毕。
\end{proof}

\begin{corollary}[Stieltjes 变换的显式线性互推：resolvent 刚性]\label{cor:pom-Lk-stieltjes-linear-rigidity}
\leanverified{paper\_pom\_lk\_stieltjes\_linear\_rigidity}
对 \(z\in\CC\setminus[0,4]\)，定义体相与边界的 Stieltjes 变换
\[
G_\nu(z):=\int_{0}^{4}\frac{1}{z-\mu}\,\dd\nu(\mu),
\qquad
G_\rho(z):=\int_{0}^{4}\frac{1}{z-\mu}\,\dd\rho(\mu).
\]
则成立显式恒等式
\[
\boxed{\;
G_\rho(z)=\frac{z-2}{2}-\frac{z(z-4)}{2}\,G_\nu(z)\; }.
\]
因此 \(G_\nu\) 与 \(G_\rho\) 互为一次变换；特别地，一者的平方根分支结构（反正弦型代数性）强制另一者同型。
\end{corollary}
\begin{proof}
由推论 \ref{cor:pom-Lk-boundary-bulk-radon-nikodym}，
\[
G_\rho(z)=\frac12\int_{0}^{4}\frac{\mu(4-\mu)}{z-\mu}\,\dd\nu(\mu)
=\frac12\int_{0}^{4}\frac{4\mu-\mu^2}{z-\mu}\,\dd\nu(\mu).
\]
利用恒等式
\[
\frac{\mu}{z-\mu}=\frac{z}{z-\mu}-1,
\qquad
\frac{\mu^2}{z-\mu}=\frac{z^2}{z-\mu}-z-\mu,
\]
得到
\[
G_\rho(z)
=\frac12\Bigl((4z-z^2)G_\nu(z)+(z-4)\underbrace{\int\dd\nu}_{=1}+\underbrace{\int\mu\,\dd\nu}_{=2}\Bigr).
\]
其中 \(\int\dd\nu=1\) 由 \(\nu\) 为概率测度得到；
而由定理 \ref{thm:pom-Lk-arcsine-law} 的 \(\theta\)-积分表示可得
\[
\int_{0}^{4}\mu\,\dd\nu(\mu)
=\frac{1}{\pi}\int_{0}^{\pi}2(1-\cos\theta)\,\dd\theta
=2.
\]
整理并注意 \(4z-z^2=-z(z-4)\) 即得所示恒等式。证毕。
\end{proof}

\begin{proposition}[体相与边界极限谱测度的 Mellin 变换闭式与有理互推]\label{prop:pom-Lk-mellin-transform-rigidity}
\leanverified{paper\_pom\_lk\_mellin\_transform\_rigidity}
定义 Mellin 变换
\[
M_\nu(s):=\int_{0}^{4}\mu^{-s}\,\dd\nu(\mu),
\qquad
M_\rho(s):=\int_{0}^{4}\mu^{-s}\,\dd\rho(\mu).
\]
则 \(M_\nu\) 的初始收敛域为 \(\Re(s)<\tfrac12\)，\(M_\rho\) 的初始收敛域为 \(\Re(s)<\tfrac32\)。
并且二者具有 Gamma 闭式（从而可作解析延拓）：
\[
M_\nu(s)=4^{-s}\,\frac{1}{\sqrt\pi}\,\frac{\Gamma(\tfrac12-s)}{\Gamma(1-s)},
\qquad
M_\rho(s)=4^{-s}\,\frac{4}{\sqrt\pi}\,\frac{\Gamma(\tfrac32-s)}{\Gamma(3-s)}.
\]
进一步，二者满足严格有理互推公式
\[
\boxed{\;
M_\rho(s)=4\,\frac{\tfrac12-s}{(1-s)(2-s)}\,M_\nu(s)\; }.
\]
\end{proposition}
\begin{proof}
由定理 \ref{thm:pom-Lk-arcsine-law} 的密度表达，作代换 \(\mu=4x\)（\(x\in(0,1)\)），得到
\[
M_\nu(s)
=\int_{0}^{1}(4x)^{-s}\,\frac{1}{\pi\sqrt{x(1-x)}}\,\dd x
=4^{-s}\,\frac{1}{\pi}B\!\Bigl(\tfrac12-s,\tfrac12\Bigr)
=4^{-s}\,\frac{1}{\sqrt\pi}\,\frac{\Gamma(\tfrac12-s)}{\Gamma(1-s)}.
\]
同理由定理 \ref{thm:pom-Lk-boundary-spectral-measure-beta32} 的密度表达，
\[
M_\rho(s)
=\int_{0}^{1}(4x)^{-s}\,\frac{8}{\pi}\,x^{1/2}(1-x)^{1/2}\,\dd x
=4^{-s}\,\frac{8}{\pi}B\!\Bigl(\tfrac32-s,\tfrac32\Bigr)
=4^{-s}\,\frac{4}{\sqrt\pi}\,\frac{\Gamma(\tfrac32-s)}{\Gamma(3-s)}.
\]
最后，由 \(\Gamma(\tfrac32-s)=(\tfrac12-s)\Gamma(\tfrac12-s)\) 与 \(\Gamma(3-s)=(2-s)(1-s)\Gamma(1-s)\) 立得互推公式。证毕。
\end{proof}


\begin{corollary}[离散 Weyl 定律：计数函数的显式反函数与统一余项界]\label{cor:pom-Lk-discrete-weyl}
定义计数函数 \(N_k(\mu):=\#\{p\le k:\ \mu_p\le \mu\}\)（\(\mu\in(0,4)\)），并令
\[
\theta(\mu):=\arccos\Bigl(1-\frac{\mu}{2}\Bigr)\in(0,\pi).
\]
则对所有 \(\mu\in(0,4)\) 有
\[
N_k(\mu)=\left\lfloor \frac{2k+1}{2\pi}\,\theta(\mu)+\frac12\right\rfloor,
\qquad
\Bigl|N_k(\mu)-\frac{2k+1}{2\pi}\theta(\mu)\Bigr|\le 1.
\]
因此 \(\nu_k\) 的分布函数在全区间上以 \(O(1/k)\) 精度逼近反正弦律的分布函数。
\end{corollary}
\begin{proof}
由于 \(\mu_p=2(1-\cos\theta_p)\) 且 \(\theta\mapsto 2(1-\cos\theta)\) 在 \((0,\pi)\) 上严格递增，
\(\mu_p\le\mu\) 等价于 \(\theta_p\le \theta(\mu)\)。
而 \(\theta_p=\frac{(2p-1)\pi}{2k+1}\)，故
\[
\theta_p\le\theta(\mu)
\quad\Longleftrightarrow\quad
2p-1\le \frac{2k+1}{\pi}\,\theta(\mu)
\quad\Longleftrightarrow\quad
p\le \frac{2k+1}{2\pi}\,\theta(\mu)+\frac12.
\]
取最大整数 \(p\) 即得显式取整式；余项界由 \(|\lfloor x\rfloor-x|\le 1\) 立得。证毕。
\end{proof}
\leanverified{paper\_pom\_lk\_discrete\_weyl}

\begin{corollary}[热核密度的闭式极限：Bessel--反正弦对偶]\label{cor:pom-Lk-heat-kernel-bessel}
定义热迹
\(
H_k(t):=\operatorname{tr}(e^{-tL_k})=\sum_{p=1}^{k}e^{-t\mu_p}
\)（\(t\ge 0\)）。
则对任意 \(t\ge 0\)，有极限
\[
\lim_{k\to\infty}\frac{1}{k}H_k(t)
=\frac{1}{\pi}\int_0^\pi e^{-2t(1-\cos\theta)}\,\dd\theta
=e^{-2t}I_0(2t),
\]
其中 \(I_0\) 为第一类修正 Bessel 函数。
并且按“边界长度”\(2k+1\) 归一化时，
\[
\lim_{k\to\infty}\frac{1}{2k+1}H_k(t)=\frac12\,e^{-2t}I_0(2t).
\]
\end{corollary}
\begin{proof}
由定理 \ref{thm:pom-Lk-arcsine-law} 取 \(f(\mu)=e^{-t\mu}\) 得第一式的积分表示。
再用 Bessel 函数的经典积分表示
\(
I_0(z)=\frac1\pi\int_0^\pi e^{z\cos\theta}\,\dd\theta
\)
并取 \(z=2t\)，得
\(
\frac{1}{\pi}\int_0^\pi e^{-2t(1-\cos\theta)}\,\dd\theta=e^{-2t}I_0(2t)
\)。
最后一式由 \(\frac{k}{2k+1}\to \frac12\) 立得。证毕。
\end{proof}

\begin{proposition}[谱 \texorpdfstring{\(\zeta\)}{zeta} 函数的主导项与半奇整数 Dirichlet 因子]\label{prop:pom-Lk-spectral-zeta-dirichlet}
对复参数 \(s\)（取 \(\Re(s)>\tfrac12\)），定义谱 \(\zeta\) 函数
\[
\zeta_k(s):=\sum_{p=1}^{k}\mu_p^{-s}.
\]
则当 \(k\to\infty\) 时有主导渐近
\[
\zeta_k(s)
=\Bigl(\frac{2k+1}{\pi}\Bigr)^{2s}\Bigl(1-2^{-2s}\Bigr)\zeta(2s)
\;+\;O(k),
\]
其中 \(\zeta\) 为 Riemann \(\zeta\) 函数。
因子 \(\bigl(1-2^{-2s}\bigr)\) 是半奇整数模态（仅出现奇频）在 \(\zeta\) 主项中所强制产生的 Dirichlet 投影因子。
\end{proposition}
\begin{proof}
记 \(n:=2k+1\) 且 \(\phi_p=\frac{(2p-1)\pi}{2n}\in(0,\pi/2)\)，则 \(\mu_p=4\sin^2\phi_p\) 且
\[
\zeta_k(s)=\sum_{p=1}^{k}(2\sin\phi_p)^{-2s}.
\]
对小角展开，\(\sin x=x+O(x^3)\) 给出 \((2\sin x)^{-2s}=(2x)^{-2s}\bigl(1+O(x^2)\bigr)\)（当 \(\Re(s)>\tfrac12\) 时该误差在求和中可积）。
将 \(2\phi_p=(2p-1)\pi/n\) 代入得
\[
(2\sin\phi_p)^{-2s}
=\Bigl(\frac{n}{\pi}\Bigr)^{2s}(2p-1)^{-2s}
\;+\;O\!\Bigl(n^{2\Re(s)-2}(2p-1)^{2-2\Re(s)}\Bigr).
\]
对 \(p\le k\) 求和，误差项为
\(
O\bigl(n^{2\Re(s)-2}\sum_{p\le k}(2p-1)^{2-2\Re(s)}\bigr)=O(n)=O(k)
\)。
因此
\[
\zeta_k(s)=\Bigl(\frac{n}{\pi}\Bigr)^{2s}\sum_{p=1}^{k}(2p-1)^{-2s}+O(k).
\]
而 \(\sum_{p=1}^{k}(2p-1)^{-2s}=(1-2^{-2s})\zeta(2s)+O(k^{1-2\Re(s)})\)；
将尾项乘以 \(n^{2\Re(s)}\) 得到 \(O(k)\) 量级，从而结论成立。证毕。
\end{proof}
\leanverified{paper\_pom\_lk\_spectral\_zeta\_dirichlet}

\begin{corollary}[归一化谱 \texorpdfstring{\(\zeta\)}{zeta} 的极限与奇频 Dirichlet 投影]\label{cor:pom-Lk-normalized-zeta-odd-projection}
对 \(\Re(s)>\tfrac12\)，定义归一化谱函数
\[
\widetilde\zeta_k(s):=\Bigl(\frac{\pi}{2k+1}\Bigr)^{2s}\zeta_k(s).
\]
则对每个 \(\Re(s)>\tfrac12\) 有极限
\[
\lim_{k\to\infty}\widetilde\zeta_k(s)
=\Bigl(1-2^{-2s}\Bigr)\zeta(2s)
=\sum_{\substack{n\ge 1\\ n\ \mathrm{odd}}}\frac{1}{n^{2s}}.
\]
\end{corollary}
\begin{proof}
将命题 \ref{prop:pom-Lk-spectral-zeta-dirichlet} 的渐近式两端乘以 \(\bigl(\pi/(2k+1)\bigr)^{2s}\) 得
\[
\widetilde\zeta_k(s)
=\Bigl(1-2^{-2s}\Bigr)\zeta(2s)
\;+\;
O(k)\Bigl(\frac{\pi}{2k+1}\Bigr)^{2s}.
\]
当 \(\Re(s)>\tfrac12\) 时，误差项为 \(O(k^{1-2\Re(s)})\to 0\)，故极限成立。
最后一式由
\(
(1-2^{-2s})\zeta(2s)=\sum_{n\ \mathrm{odd}} n^{-2s}
\)
的恒等式给出。证毕。
\end{proof}

\begin{corollary}[连续极限的严格匹配：ND 行列式极限与 \texorpdfstring{\(t/4\)}{t/4} 的唯一尺度]\label{cor:pom-Lk-continuum-ND-det-cosh}
令连续算子 \(\mathcal L:=-\frac{\dd^2}{\dd x^2}\) 作用于区间 \([0,1]\)，取 Neumann--Dirichlet 边界条件 \(u'(0)=0,\ u(1)=0\)。
其谱为 \(\{((2p-1)\pi/2)^2\}_{p\ge 1}\)，且其加质量行列式满足
\(
\det(\mathcal L+s)=\cosh(\sqrt{s})
\)（\(\zeta\)-正规化意义）。
\par
另一方面，离散谱在小模态上满足
\[
\mu_p=\Bigl(\frac{(2p-1)\pi}{2k+1}\Bigr)^2+O(k^{-4})\qquad(p\ \text{固定},\ k\to\infty),
\]
因此缩放算子 \(\widetilde L_k:=\frac{(2k+1)^2}{4}L_k\) 在谱意义下趋近于 \(\mathcal L\)。
更强地，对任意固定 \(s\ge 0\) 有行列式极限
\[
\lim_{k\to\infty}\det\!\Bigl(L_k+\frac{4s}{(2k+1)^2}I\Bigr)=\cosh(\sqrt{s}).
\]
从而尺度因子 \(t/4\)（即以 \(4s/(2k+1)^2\) 加质量）并非可选规范，而由半奇整数模态的连续极限匹配结构性强制给出。
\end{corollary}
\begin{proof}
取 \(n:=2k+1\)。由推论 \ref{cor:pom-Lk-shifted-det-free-energy}，
\[
\det\!\Bigl(L_k+\frac{4s}{n^2}I\Bigr)
=\frac{T_{n}\!\bigl(\sqrt{1+s/n^2}\bigr)}{\sqrt{1+s/n^2}}.
\]
当 \(x:=\sqrt{1+s/n^2}\ge 1\) 时，有 \(T_n(x)=\cosh(n\,\mathrm{arcosh}(x))\)。
并且 \(n\,\mathrm{arcosh}(\sqrt{1+s/n^2})=\sqrt{s}+O(n^{-2})\)（由 \(\mathrm{arcosh}(1+u)=\sqrt{2u}+O(u^{3/2})\) 及 \(u=s/(2n^2)+O(n^{-4})\) 得到）。
于是
\[
\det\!\Bigl(L_k+\frac{4s}{n^2}I\Bigr)
=\frac{\cosh(\sqrt{s}+O(n^{-2}))}{1+O(n^{-2})}
\longrightarrow \cosh(\sqrt{s}).
\]
最后关于“小模态匹配”的特征值渐近由 \(\sin x=x+O(x^3)\) 直接给出。证毕。
\end{proof}

\begin{corollary}[ND 行列式极限的局部一致收敛：谱不变量可逐阶交换极限]\label{cor:pom-Lk-continuum-ND-det-cosh-locally-uniform}
沿用推论 \ref{cor:pom-Lk-continuum-ND-det-cosh} 的记号，并设
\[
D_k(s):=\det\!\Bigl(L_k+\frac{4s}{(2k+1)^2}I\Bigr).
\]
则对任意紧集 \(K\subset\CC\) 有局部一致收敛
\[
\sup_{s\in K}\bigl|D_k(s)-\cosh(\sqrt{s})\bigr|\longrightarrow 0,
\]
其中 \(\cosh(\sqrt{s})\) 视为整函数
\(
\sum_{n\ge 0}s^n/(2n)!
\)。
因此可逐阶交换导数与极限：对任意 \(n\ge 0\) 成立
\[
\lim_{k\to\infty}D_k^{(n)}(0)=\frac{n!}{(2n)!}.
\]
\end{corollary}
\begin{proof}
令 \(n:=2k+1\)。由推论 \ref{cor:pom-Lk-shifted-det-free-energy}，
\[
D_k(s)=\frac{T_{n}\!\bigl(\sqrt{1+s/n^2}\bigr)}{\sqrt{1+s/n^2}}.
\]
固定紧集 \(K\subset\CC\)，记 \(R:=\max_{s\in K}|s|\)。
当 \(n\) 充分大时 \(|s|/n^2\le R/n^2\) 足够小，从而
\(\sqrt{1+s/n^2}=1+O_K(n^{-2})\) 一致成立。
并且由 \(\mathrm{arsinh}(z)=z+O(|z|^3)\) 得
\[
n\,\mathrm{arsinh}\!\Bigl(\frac{\sqrt{s}}{n}\Bigr)=\sqrt{s}+O_K(n^{-2}),
\]
其中误差界仅依赖于 \(|s|\)，故对 \(K\) 一致。
结合恒等式 \(T_n(\sqrt{1+z^2})=\cosh(n\,\mathrm{arsinh}(z))\)（先在 \(z\in\RR_{\ge 0}\) 上成立，再由解析延拓推广），得到
\[
T_n(\sqrt{1+s/n^2})=\cosh(\sqrt{s}+O_K(n^{-2})).
\]
由于 \(\cosh\) 在有界集上一致 Lipschitz，且分母 \(\sqrt{1+s/n^2}=1+O_K(n^{-2})\)，可得
\(
\sup_{s\in K}|D_k(s)-\cosh(\sqrt{s})|\to 0
\)。
最后，局部一致收敛由 Weierstrass 定理推出导数可逐阶交换，从而
\(
\lim_{k\to\infty}D_k^{(n)}(0)=\bigl(\cosh(\sqrt{s})\bigr)^{(n)}|_{s=0}=n!/(2n)!
\)。
证毕。
\end{proof}

\begin{remark}[有限窗审计：经验谱律、热核与 ND 行列式极限]\label{rem:pom-fence-green-kernel-spectral-limit-laws-audit}
脚本在有限窗上核验：经验谱分布与反正弦律的分布函数误差为 \(O(1/k)\)（推论 \ref{cor:pom-Lk-discrete-weyl}），
热核密度 \(\frac{1}{k}H_k(t)\) 对 \(e^{-2t}I_0(2t)\) 的收敛（推论 \ref{cor:pom-Lk-heat-kernel-bessel}），以及
\(\det\bigl(L_k+\frac{4s}{(2k+1)^2}I\bigr)\) 对 \(\cosh(\sqrt{s})\) 的收敛（推论 \ref{cor:pom-Lk-continuum-ND-det-cosh}）。
审计表如下：
\begin{table}[H]
\centering
\scriptsize
\setlength{\tabcolsep}{6pt}
\caption{混合边界离散 Laplacian $L_k=K_k^{-1}$ 的有限窗极限律审计：经验谱分布对反正弦律的 Kolmogorov 距离（定理~\ref{thm:pom-Lk-arcsine-law}），离散 Weyl 计数余项上界（推论~\ref{cor:pom-Lk-discrete-weyl}），热核密度对 $e^{-2t}I_0(2t)$ 的收敛（推论~\ref{cor:pom-Lk-heat-kernel-bessel}，取 $t=1$），以及 ND 行列式极限 $\cosh(\sqrt{s})$ 的收敛（推论~\ref{cor:pom-Lk-continuum-ND-det-cosh}，取 $s=1$）。}
\label{tab:fence_green_kernel_spectral_limit_laws_audit}
\begin{tabular}{r r r r r}
\toprule
$k$ & $\|F_k-F\|_{\infty}$ & $\sup|R_k|$ & $|H_k(1)/k- e^{-2}I_0(2)|$ & $|\det-\cosh(1)|$\\
\midrule
50 & 1.98e-02 & 5.00e-01 & 2.90e-03 & 9.48e-05\\
100 & 9.95e-03 & 5.00e-01 & 1.45e-03 & 2.39e-05\\
200 & 4.99e-03 & 5.00e-01 & 7.25e-04 & 6.02e-06\\
400 & 2.50e-03 & 5.00e-01 & 3.63e-04 & 1.51e-06\\
\bottomrule
\end{tabular}
\end{table}

\end{remark}
